% Vietnamese translation for OI Public Library
% Source-PDF: 模板 TEMPLATES/chota1219.pdf
\documentclass[11pt]{article}
\usepackage{oipl-vn}
\usepackage{fvextra}
\setmonofont{Unifont}[Scale=MatchLowercase]
\DefineVerbatimEnvironment{Code}{Verbatim}{fontsize=\scriptsize,breaklines=true,breakanywhere=true}
\DefineVerbatimEnvironment{NoteBlock}{Verbatim}{fontsize=\small,breaklines=true,breakanywhere=true}
\title{Sổ tay mẫu mã chota1219}
\author{Bản dịch tiếng Việt}
\date{}
\begin{document}
\maketitle
\section{1 Mẫu khung chương trình}
\begin{Code}
1 import static java.lang.Math.*;
2 import static java.util.Arrays.*;
3 import java.io.*;
4 import java.util.*;
5
6 public class Main {
7 static boolean LOCAL = System.getSecurityManager() == null;
8 Scanner sc = new Scanner(System.in);
9
10 void run() {
11 }
12
13 void debug(Object...os) {
14 System.err.println(deepToString(os));
15 }
16
17 public static void main(String[] args) {
18 if (LOCAL) {
19 try {
20 System.setIn(new FileInputStream("in.txt"));
21 } catch (Throwable e) {
22 LOCAL = false;
23 }
24 }
25 new Main().run();
26 }
27 }
\end{Code}
\section{2 Cấu trúc dữ liệu}
\subsection{2.1 BIT}
\begin{Code}
1 class BIT {
2 int[] vs;
3 BIT(int n) {
4 vs = new int[n + 1];
5 }
6 void add(int k, int a) {
7 for (int i = k + 1; i < vs.length; i += i & -i) {
8 vs[i] += a;
9 }
10 }
11 int sum(int s, int t) {
12 if (s > 0) return sum(0, t) - sum(0, s);
13 int res = 0;
14 for (int i = t; i > 0; i -= i & -i) {
15 res += vs[i];
16 }
17 return res;
18 }
19 //[0,i] の和が k より大きくなる最小の i を求める
20 int get(int k) {
21 int p = Integer.highestOneBit(vs.length - 1);
22 for (int q = p; q > 0; q >>= 1, p |= q) {
23 if (p >= vs.length || k < vs[p]) p ^= q;
24 else k -= vs[p];
25 }
26 return p;
27 }
28 }
\end{Code}
\subsection{2.2 RMQ}
\begin{Code}
1 class RMQ {
2 int[] vs;
3 int[][] min;
4 RMQ(int[] vs) {
5 int n = vs.length, m = log2(n) + 1;
6 this.vs = vs;
7 min = new int[m][n];
8 for (int i = 0; i < n; i++) min[0][i] = i;
9 for (int i = 1, k = 1; i < m; i++, k <<= 1) {
10 for (int j = 0; j + k < n; j++) {
11 min[i][j] = vs[min[i - 1][j]] <= vs[min[i - 1][j + k]] ?
12 min[i - 1][j] : min[i - 1][j + k];
13 }
14 }
15 }
16 //最小値が複数存在する場合は一番最初のを返す
17 int query(int from, int to) {
18 int k = log2(to - from);
19 return vs[min[k][from]] <= vs[min[k][to - (1 << k)]] ?
20 min[k][from] : min[k][to - (1 << k)];
21 }
22 int log2(int b) {
23 return 31 - Integer.numberOfLeadingZeros(b);
24 }
25 }
\end{Code}
\subsection{2.3 Cập nhật đoạn}
\begin{Code}
1 class Intervals {
2 TreeMap<Integer, Integer> map = new TreeMap<Integer, Integer>();
3 Intervals() {
4 map.put(Integer.MIN_VALUE, -1);
5 map.put(Integer.MAX_VALUE, -1);
6 }
7 void paint(int s, int t, int c) {
8 int p = get(t);
9 map.subMap(s, t).clear();
10 map.put(s, c);
11 map.put(t, p);
12 }
13 int get(int k) {
14 return map.floorEntry(k).getValue();
15 }
16 }
\end{Code}
\subsection{2.4 Treap}
\begin{NoteBlock}
Cài đặt không gây tác dụng phụ.
\end{NoteBlock}
\begin{Code}
1 class T {
2 final int key, val;
3 final double p;
4 final T left, right;
5 T(int key, int val, double p, T left, T right) {
6 this.key = key;
7 this.val = val;
8 this.p = p;
9 this.left = left;
10 this.right = right;
11 }
12 T change(T left, T right) {
13 return new T(key, val, p, left, right);
14 }
15 T normal() {
16 if (left != null && left.p < p && (right == null || left.p < right.p)) {
17 return left.change(left.left, change(left.right, right));
18 } else if (right != null && right.p < p) {
19 return right.change(change(left, right.left), right.right);
20 }
21 return this;
22 }
23 }
24 T put(T t, int key, int val) {
25 if (t == null) return new T(key, val, random(), null, null);
26 if (key < t.key) return t.change(put(t.left, key, val), t.right).normal();
27 if (key > t.key) return t.change(t.left, put(t.right, key, val)).normal();
28 return new T(key, val, t.p, t.left, t.right);
29 }
30 T remove(T t, int key) {
31 if (t == null) return null;
32 if (key < t.key) return t.change(remove(t.left, key), t.right);
33 if (key > t.key) return t.change(t.left, remove(t.right, key));
34 return merge(t.left, t.right);
35 }
36 T merge(T t1, T t2) {
37 if (t1 == null) return t2;
38 if (t2 == null) return t1;
39 if (t1.p < t2.p) return t1.change(t1.left, merge(t1.right, t2));
40 return t2.change(merge(t1, t2.left), t2.right);
41 }
\end{Code}
\subsection{2.5 Cây đoạn}
\begin{NoteBlock}
Cài đặt nhanh bằng mảng và vòng lặp.
parent(i) = i/2
left(i) = 2i
size(i) = N/highest(i)
s(i) = size(i)i - N
\end{NoteBlock}
\begin{Code}
1 class Seg {
2 int N;
3 Seg(int n) {
4 N = Integer.highestOneBit(n) << 1;
5 }
6 //点を含む区間に対する処理
7 void query(int k) {
8 for (int i = N + k; i > 0; i >>= 1) {
9 //...
10 }
11 }
12 //カバーする区間に対する処理
13 void query(int s, int t) {
14 while (0 < s && s + (s & -s) <= t) {
15 int i = (N + s) / (s & -s);
16 //...
17 s += s & -s;
18 }
19 while (s < t) {
20 int i = (N + t) / (t & -t) - 1;
21 //...
22 t -= t & -t;
23 }
24 }
25 }
\end{Code}
\section{3 Đồ thị}
\subsection{3.1 Tập công thức}
\subsubsection{3.1.1 Định lý đa diện Euler}
\begin{NoteBlock}
Ghi chú/công thức: V ,   E,   F   V - E + F = 2
\end{NoteBlock}
\subsubsection{3.1.2 Định lý max-flow min-cut}
\begin{NoteBlock}
Giá trị luồng s-t lớn nhất bằng dung lượng lát cắt s-t nhỏ nhất.
Nếu X là tập các đỉnh đi được từ s trong đồ thị dư Gf, tập cạnh của min-cut trong G là delta^+_G(X).
G(X)
\end{NoteBlock}
\subsubsection{3.1.3 Các hệ thức về matching}
\begin{NoteBlock}
|matching lớn nhất| <= |vertex cover nhỏ nhất|; với đồ thị hai phía thì có dấu bằng.
|tập độc lập lớn nhất| + |vertex cover nhỏ nhất| = V.
Với đồ thị không có đỉnh cô lập: |matching lớn nhất| + |edge cover nhỏ nhất| = V.
Với đồ thị hai phía: |edge cover nhỏ nhất| = |tập độc lập lớn nhất|.
X subset V(G) là vertex cover của G khi và chỉ khi V(G) \ X là tập độc lập của G.
\end{NoteBlock}
\subsubsection{3.1.4 Định lý cây ma trận}
\begin{NoteBlock}
Ghi chú/công thức: G   (G   - G  )
\end{NoteBlock}
\subsubsection{3.1.5 Số cây theo dãy bậc}
\begin{NoteBlock}
Trong K_n, số cây khung sao cho đỉnh i có bậc d_i là:
(n - 2)!
(d1 - 1)!(d2 - 1)! . . . (dn - 1)!
\end{NoteBlock}
\subsubsection{3.1.6 Công thức Kasteleyn}
\begin{NoteBlock}
Với đồ thị có thể định hướng mọi chu trình độ dài chẵn theo kiểu lẻ, định thức ma trận kề của nó (cạnh ngược hướng lấy -1)
bằng bình phương số matching hoàn hảo.
Đồ thị phẳng luôn có một định hướng như vậy.
Với đồ thị hai phía, định thức ma trận kề hai phía bằng số matching hoàn hảo.
\end{NoteBlock}
\subsubsection{3.1.7 Ma trận Tutte}
\begin{NoteBlock}
Ma trận Tutte là ma trận V x V thu được bằng cách gán hướng tùy ý cho các cạnh của G để được đồ thị có hướng G' .
Với cạnh (v, u),
Ghi chú/công thức: xvu  tvu = xvu, tuv = -xvu
Với đồ thị tổng quát, hạng của ma trận Tutte bằng hai lần kích thước matching.
Lấy x ngẫu nhiên cho nghiệm đúng với xác suất đủ cao.
\end{NoteBlock}
\subsection{3.2 Sắp xếp tô pô}
\begin{NoteBlock}
Thứ tự tô pô là thứ tự thỏa: nếu có cạnh từ v[i] đến v[j] thì i < j.
Nếu đánh số lúc thoát khỏi DFS, ta nhận được thứ tự tô pô đảo.
\end{NoteBlock}
\begin{Code}
1 V[] topologicalSort(V[] vs) {
2 n = vs.length;
3 us = new V[n];
4 for (V v : vs) {
5 if (v.state == 0 && !dfs(v)) return null;
6 }
7 return us;
8 }
9 boolean dfs(V v) {
10 v.state = 1;
11 for (V u : v) {
12 if (u.state == 1 || u.state == 0 && !dfs(u)) return false;
13 }
14 us[--n] = v;
15 v.state = 2;
16 return true;
17 }
\end{Code}
\subsection{3.3 Phân rã thành phần liên thông mạnh}
\begin{NoteBlock}
Trước hết tìm thứ tự tô pô (bỏ qua chu trình), rồi DFS theo cạnh đảo theo thứ tự đó.
Các cây tìm kiếm cực đại thu được khi đó là các thành phần liên thông mạnh.
Trả về số SCC là k; comp chứa thứ tự tô pô trong [0,k).
\end{NoteBlock}
\begin{Code}
1 int scc(V[] vs) {
2 n = vs.length;
3 us = new V[n];
4 for (V v : vs) if (!v.visit) dfs(v);
5 for (V v : vs) v.visit = false;
6 for (V u : us) if (!u.visit) dfsrev(u, n++);
7 return n;
8 }
9 void dfs(V v) {
10 v.visit = true;
11 for (V u : v.fs) if (!u.visit) dfs(u);
12 us[--n] = v;
13 }
14 void dfsrev(V v, int k) {
15 v.visit = true;
16 for (V u : v.rs) if (!u.visit) dfsrev(u, k);
17 v.comp = k;
18 }
\end{Code}
\subsection{3.4 Đỉnh khớp và cầu}
\begin{NoteBlock}
Chọn gốc tùy ý, chạy DFS và gán số num khi đi vào đỉnh.
Tính low là giá trị nhỏ nhất trong các mục sau:
- num[v].
- num[u] của đỉnh u đã thăm, kề với v, không phải cạnh vừa đi tới.
- low[u] của con u của v.
Khi đó:
Gốc: r là đỉnh khớp khi và chỉ khi r có ít nhất hai con.
Không phải gốc: v là đỉnh khớp khi và chỉ khi tồn tại con u sao cho num[v] <= low[u].
(v,u) là cầu khi và chỉ khi num[v] < low[u].
count chứa số phần tách ra khi bỏ đỉnh đó; hàm trả về tập cầu.
\end{NoteBlock}
\begin{Code}
1 ArrayList<E> connection(V[] vs) {
2 bridge = new ArrayList<E>();
3 for (V v : vs) if (v.num < 0) {
4 dfs(v, 0);
5 if (v.count > 0) v.count--;
6 }
7 return bridge;
8 }
9 int dfs(V v, int c) {
10 v.num = c;
11 int low = c;
12 boolean rev = false;
13 for (V u : v) {
14 if (u.num < 0) {
15 int t = dfs(u, c + 1);
16 low = min(low, t);
17 if (v.num <= t) v.count++;
18 if (v.num < t) bridge.add(new E(v, u));
19 } else if (u.num != v.num - 1 || rev) low = min(low, u.num);
20 else rev = true;
21 }
22 return low;
23 }
\end{Code}
\subsection{3.5 Chu trình có độ dài trung bình nhỏ nhất}
\begin{NoteBlock}
Cho phép cạnh âm.
Bất kỳ chu trình nào nằm trên đường đi truy vết bằng prev đều là chu trình có trung bình nhỏ nhất.
O(V E)
\end{NoteBlock}
\begin{Code}
1 double mmc(int[] ss, int[] ts, double[] cs, int n) {
2 int m = ss.length;
3 double[][] dp = new double[n + 1][n];
4 for (int i = 0; i < n; i++) {
5 fill(dp[i + 1], INF);
6 for (int j = 0; j < m; j++) {
7 dp[i + 1][ts[j]] = min(dp[i + 1][ts[j]], dp[i][ss[j]] + cs[j]);
8 }
9 }
10 double res = INF;
11 for (int i = 0; i < n; i++) if (dp[n][i] < INF) {
12 double max = -INF;
13 for (int j = 0; j < n; j++) {
14 max = max(max, (dp[n][i] - dp[j][i]) / (n - j));
15 }
16 res = min(res, max);
17 }
18 return res;
19 }
\end{Code}
\subsection{3.6 Cây có hướng nhỏ nhất với gốc cho trước}
\begin{NoteBlock}
Tham lam chọn cạnh vào nhỏ nhất cho mỗi đỉnh khác gốc; nếu sinh chu trình thì co lại và xử lý đệ quy.
Muốn khôi phục cây có hướng, hãy lưu cạnh gốc khi co.
Nếu chưa chỉ định gốc, trừ một giá trị đủ lớn khỏi mọi cạnh, rồi thêm cạnh chi phí 0 từ gốc đến mọi đỉnh khác
là đủ.
Muốn tìm rừng có hướng lớn nhất, đảo dấu trọng số cạnh rồi thêm cạnh chi phí 0 từ gốc đến mọi đỉnh khác.
O(V E)
\end{NoteBlock}
\begin{Code}
1 int arborescence(V[] vs, V r) {
2 int res = 0;
3 for (V v : vs) for (E e : v.es) e.to.min = min(e.to.min, e.cost);
4 for (V v : vs) if (v != r) {
5 if (v.min == INF) return INF;
6 res += v.min;
7 }
8 for (V v : vs) for (E e : v.es) if (e.to != r) {
9 e.cost -= e.to.min;
10 if (e.cost == 0) {
11 v.fs.add(e.to);
12 e.to.rs.add(v);
13 }
14 }
15 int m = scc(vs);
16 if (m == vs.length) return res;
17 V[] us = new V[m];
18 for (int i = 0; i < m; i++) us[i] = new V();
19 for (V v : vs) for (E e : v.es) {
20 if (v.comp != e.to.comp) us[v.comp].add(us[e.to.comp], e.cost);
21 }
22 return min(INF, res + arborescence(us, us[r.comp]));
23 }
\end{Code}
\subsection{3.7 Cây Steiner}
\begin{NoteBlock}
Với đồ thị vô hướng không âm, tìm chi phí cây nhỏ nhất chứa mọi terminal trong ts.
dp[S][v] := chi phí cây Steiner nhỏ nhất chứa S subset T và v in V.
O(3t
n + 2t
n2
+ n3
)
\end{NoteBlock}
\begin{Code}
1 int steiner(int[][] g, int[] ts) {
2 int n = g.length, m = ts.length;
3 if (m < 2) return 0;
4 int[][] dp = new int[1 << m][n];
5 for (int k = 0; k < n; k++) {
6 for (int i = 0; i < n; i++) {
7 for (int j = 0; j < n; j++) {
8 g[i][j] = min(g[i][j], g[i][k] + g[k][j]);
9 }
10 }
11 }
12 for (int i = 0; i < m; i++) {
13 for (int j = 0; j < n; j++) {
14 dp[1 << i][j] = g[ts[i]][j];
15 }
16 }
17 for (int i = 1; i < 1 << m; i++) if (((i - 1) & i) != 0) {
18 for (int j = 0; j < n; j++) {
19 dp[i][j] = INF;
20 for (int k = (i - 1) & i; k > 0; k = (k - 1) & i) {
21 dp[i][j] = min(dp[i][j], dp[k][j] + dp[i ^ k][j]);
22 }
23 }
24 for (int j = 0; j < n; j++) {
25 for (int k = 0; k < n; k++) {
26 dp[i][j] = min(dp[i][j], dp[i][k] + g[k][j]);
27 }
28 }
29 }
30 return dp[(1 << m) - 1][ts[0]];
31 }
\end{Code}
\subsection{3.8 Dinic}
\begin{Code}
1 int dinic(V s, V t) {
2 int flow = 0;
3 for (int p = 1; ; p++) {
4 Queue<V> que = new LinkedList<V>();
5 s.level = 0;
6 s.p = p;
7 que.offer(s);
8 while (!que.isEmpty()) {
9 V v = que.poll();
10 v.iter = v.es.size() - 1;
11 for (E e : v.es) if (e.to.p < p && e.cap > 0) {
12 e.to.level = v.level + 1;
13 e.to.p = p;
14 que.offer(e.to);
15 }
16 }
17 if (t.p < p) return flow;
18 for (int f; (f = dfs(s, t, INF)) > 0; ) flow += f;
19 }
20 }
21 int dfs(V v, V t, int f) {
22 if (v == t) return f;
23 for (; v.iter >= 0; v.iter--) {
24 E e = v.es.get(v.iter);
25 if (v.level < e.to.level && e.cap > 0) {
26 int d = dfs(e.to, t, min(f, e.cap));
27 if (d > 0) {
28 e.cap -= d;
29 e.rev.cap += d;
30 return d;
31 }
32 }
33 }
34 return 0;
35 }
36 class V {
37 ArrayList<E> es = new ArrayList<E>();
38 int level, p, iter;
39 void add(V to, int cap) {
40 E e = new E(to, cap), rev = new E(this, 0);
41 e.rev = rev; rev.rev = e;
42 es.add(e); to.es.add(rev);
43 }
44 }
45 class E {
46 V to;
47 E rev;
48 int cap;
49 E(V to, int cap) {
50 this.to = to;
51 this.cap = cap;
52 }
53 }
\end{Code}
\subsection{3.9 Global min-cut nhỏ nhất}
\begin{NoteBlock}
Trọng số phải không âm.
O(n3
)
\end{NoteBlock}
\begin{Code}
1 int minCut(int[][] c) {
2 int n = c.length, cut = INF;
3 int[] id = new int[n], b = new int[n];
4 for (int i = 0; i < n; i++) id[i] = i;
5 for ( ; n > 1; n--) {
6 fill(b, 0);
7 for (int i = 0; i + 1 < n; i++) {
8 int p = i + 1;
9 for (int j = i + 1; j < n; j++) {
10 b[id[j]] += c[id[i]][id[j]];
11 if (b[id[p]] < b[id[j]]) p = j;
12 }
13 swap(id, i + 1, p);
14 }
15 cut = min(cut, b[id[n - 1]]);
16 for (int i = 0; i < n - 2; i++) {
17 c[id[i]][id[n - 2]] += c[id[i]][id[n - 1]];
18 c[id[n - 2]][id[i]] += c[id[n - 1]][id[i]];
19 }
20 }
21 return cut;
22 }
\end{Code}
\subsection{3.10 Luồng chi phí nhỏ nhất}
\begin{NoteBlock}
Đừng quên truyền vào toàn bộ các đỉnh.
\end{NoteBlock}
\begin{Code}
1 int minCostFlow(V[] vs, V s, V t, int flow) {
2 int res = 0;
3 while (flow > 0) {
4 for (V v : vs) v.min = INF;
5 PriorityQueue<E> que = new PriorityQueue<E>();
6 s.min = 0;
7 que.offer(new E(s, 0, 0));
8 while (!que.isEmpty()) {
9 E crt = que.poll();
10 if (crt.cost == crt.to.min) {
11 for (E e : crt.to.es) {
12 int tmp = crt.cost + e.cost + crt.to.h - e.to.h;
13 if (e.cap > 0 && e.to.min > tmp) {
14 e.to.min = tmp;
15 e.to.prev = e;
16 que.offer(new E(e.to, 0, e.to.min));
17 }
18 }
19 }
20 }
21 if (t.min == INF) return -1;
22 int d = flow;
23 for (E e = t.prev; e != null; e = e.rev.to.prev) {
24 d = min(d, e.cap);
25 }
26 for (E e = t.prev; e != null; e = e.rev.to.prev) {
27 res += d * e.cost;
28 e.cap -= d;
29 e.rev.cap += d;
30 }
31 flow -= d;
32 for (V v : vs) v.h += v.min;
33 }
34 return res;
35 }
36 class V {
37 ArrayList<E> es = new ArrayList<E>();
38 E prev;
39 int min, h;
40 void add(V to, int cap, int cost) {
41 E e = new E(to, cap, cost), rev = new E(this, 0, -cost);
42 e.rev = rev; rev.rev = e;
43 es.add(e); to.es.add(rev);
44 }
45 }
46 class E implements Comparable<E> {
47 V to;
48 E rev;
49 int cap, cost;
50 E(V to, int cap, int cost) {
51 this.to = to;
52 this.cap = cap;
53 this.cost = cost;
54 }
55 public int compareTo(E o) {
56 return cost - o.cost;
57 }
58 }
\end{Code}
\subsection{3.11 Matching hai phía}
\begin{NoteBlock}
Vertex cover nhỏ nhất gồm các đỉnh không đi tới được ở phía trái và các đỉnh đi tới được ở phía phải.
O(V E)
\end{NoteBlock}
\begin{Code}
1 int bipartiteMatching(V[] vs) {
2 int match = 0;
3 for (V v : vs) if (v.pair == null) {
4 for (V u : vs) u.used = false;
5 if (dfs(v)) match++;
6 }
7 return match;
8 }
9 boolean dfs(V v) {
10 v.used = true;
11 for (V u : v) {
12 V w = u.pair;
13 if (w == null || !w.used && dfs(w)) {
14 v.pair = u;
15 u.pair = v;
16 return true;
17 }
18 }
19 return false;
20 }
21 class V extends ArrayList<V> {
22 V pair;
23 boolean used;
24 void connect(V v) {
25 add(v);
26 v.add(this);
27 }
28 }
\end{Code}
\begin{NoteBlock}
O(
√
V E)
\end{NoteBlock}
\begin{Code}
1 int hopcroftKarp(V[] vs) {
2 for (int match = 0;;) {
3 Queue<V> que = new LinkedList<V>();
4 for (V v : vs) v.level = -1;
5 for (V v : vs) if (v.pair == null) {
6 v.level = 0;
7 que.offer(v);
8 }
9 while (!que.isEmpty()) {
10 V v = que.poll();
11 for (V u : v) {
12 V w = u.pair;
13 if (w != null && w.level < 0) {
14 w.level = v.level + 1;
15 que.offer(w);
16 }
17 }
18 }
19 for (V v : vs) v.used = false;
20 int d = 0;
21 for (V v : vs) if (v.pair == null && dfs(v)) d++;
22 if (d == 0) return match;
23 match += d;
24 }
25 }
26 boolean dfs(V v) {
27 v.used = true;
28 for (V u : v) {
29 V w = u.pair;
30 if (w == null || !w.used && v.level < w.level && dfs(w)) {
31 v.pair = u;
32 u.pair = v;
33 return true;
34 }
35 }
36 return false;
37 }
\end{Code}
\subsection{3.12 Matching ổn định}
\begin{NoteBlock}
Cách tìm stable matching tối ưu cho phía nam.
- Chọn một nam và cho anh ta cầu hôn người nữ ưu tiên cao nhất trong số chưa từng cầu hôn.
- Nếu người nữ thích nam này hơn bạn ghép hiện tại, hoặc chưa có bạn ghép, tạo cặp mới với nam này,
và hủy cặp cũ.
- Người nam bị từ chối lập tức cầu hôn người nữ khác.
O(n2
)
Trả về res[i] là bạn ghép của nam i.
orderM[i][j] là người nữ mà nam i thích thứ j.
preferW[i][j] là thứ hạng ưu tiên của nữ i đối với nam j.
\end{NoteBlock}
\begin{Code}
1 int[] stableMatching(int[][] orderM, int[][] preferW) {
2 int n = orderM.length;
3 int[] pairM = new int[n], pairW = new int[n], p = new int[n];
4 fill(pairM, -1);
5 fill(pairW, -1);
6 for (int i = 0; i < n; i++) {
7 while (pairM[i] < 0) {
8 int w = orderM[i][p[i]++], m = pairW[w];
9 if (m == -1) {
10 pairM[i] = w;
11 pairW[w] = i;
12 } else if (preferW[w][i] < preferW[w][m]) {
13 pairM[m] = -1;
14 pairM[i] = w;
15 pairW[w] = i;
16 i = m;
17 }
18 }
19 }
20 return pairM;
21 }
\end{Code}
\subsection{3.13 Chu trình Euler}
\begin{NoteBlock}
Đừng quên kiểm tra liên thông.
Với đồ thị vô hướng, cũng cần kiểm tra rev.
\end{NoteBlock}
\begin{Code}
1 V[] eulerianWalk(V v) {
2 Stack<V> res = new Stack<V>(), stack = new Stack<V>();
3 stack.push(v);
4 while (!stack.isEmpty()) {
5 v = stack.pop();
6 while (v.p < v.size()) {
7 stack.push(v);
8 v = v.get(v.p++);
9 }
10 res.push(v);
11 }
12 return res.toArray(new V[0]);
13 }
\end{Code}
\subsection{3.14 LCA}
\begin{NoteBlock}
Gọi e[2*n-1] là thứ tự các nút được thăm khi DFS từ gốc, v[2*n-1] là độ sâu tương ứng, và
a[n] là vị trí xuất hiện đầu tiên trong e của mỗi nút; đáp án là e[RMQ_v([a[u], a[v]])].
\end{NoteBlock}
\subsection{3.15 2SAT}
\begin{NoteBlock}
Với mỗi mệnh đề (y or z), tạo đồ thị có hai cạnh not y -> z và not z -> y.
Phân rã SCC đồ thị này; việc một SCC chứa đồng thời x và not x tương đương với công thức gốc không thỏa được
(unsatisfiable).
SCC chứa x đứng sau SCC chứa not x trong thứ tự tô pô khi và chỉ khi x là true.
\end{NoteBlock}
\section{4 Hình học phẳng}
\subsection{4.1 Giao điểm và các phép tính liên quan}
\begin{Code}
1 //線分と点の距離
2 double disSP(P p1, P p2, P q) {
3 if (p2.sub(p1).dot(q.sub(p1)) < EPS) return q.sub(p1).abs();
4 if (p1.sub(p2).dot(q.sub(p2)) < EPS) return q.sub(p2).abs();
5 return disLP(p1, p2, q);
6 }
7 //直線と点の距離
8 double disLP(P p1, P p2, P q) {
9 return abs(p2.sub(p1).det(q.sub(p1))) / p2.sub(p1).abs();
10 }
11 //線分と線分の交差判定
12 boolean crsSS(P p1, P p2, P q1, P q2) {
13 if (max(p1.x, p2.x) + EPS < min(q1.x, q2.x)) return false;
14 if (max(q1.x, q2.x) + EPS < min(p1.x, p2.x)) return false;
15 if (max(p1.y, p2.y) + EPS < min(q1.y, q2.y)) return false;
16 if (max(q1.y, q2.y) + EPS < min(p1.y, p2.y)) return false;
17 return sig(p2.sub(p1).det(q1.sub(p1))) * sig(p2.sub(p1).det(q2.sub(p1))) < EPS
18 && sig(q2.sub(q1).det(p1.sub(q1))) * sig(q2.sub(q1).det(p2.sub(q1))) < EPS;
19 }
20 //円と線分の交差判定
21 boolean crsCS(P c, double r, P p1, P p2) {
22 return disSP(p1, p2, c) < r + EPS &&
23 (r < c.sub(p1).abs() + EPS || r < c.sub(p2).abs() + EPS);
24 }
25 //円と円の交差判定
26 boolean crsCC(P c1, double r1, P c2, double r2) {
27 double dis = c1.sub(c2).abs();
28 return dis < r1 + r2 + EPS && abs(r1 - r2) < dis + EPS;
29 }
30 //四点の同一円周上判定 (一直線上は true)
31 boolean onCir(P p1, P p2, P p3, P p4) {
32 if (abs(p2.sub(p1).det(p3.sub(p1))) < EPS) return true;
33 P c = ccenter(p1, p2, p3);
34 return abs(c.sub(p1).abs2() - c.sub(p4).abs2()) < EPS;
35 }
36 //点から直線に下ろした垂線の足
37 P proj(P p1, P p2, P q) {
38 return p1.add(p2.sub(p1).mul(p2.sub(p1).dot(q.sub(p1)) / p2.sub(p1).abs2()));
39 }
40 //直線と直線の交点
41 P isLL(P p1, P p2, P q1, P q2) {
42 double d = q2.sub(q1).det(p2.sub(p1));
43 if (abs(d) < EPS) return null;
44 return p1.add(p2.sub(p1).mul(q2.sub(q1).det(q1.sub(p1)) / d));
45 }
46 //直線と円の交点 (p1 に近い順)
47 P[] isCL(P c, double r, P p1, P p2) {
48 double x = p1.sub(c).dot(p2.sub(p1));
49 double y = p2.sub(p1).abs2();
50 double d = x * x - y * (p1.sub(c).abs2() - r * r);
51 if (d < -EPS) return new P[0];
52 if (d < 0) d = 0;
53 P q1 = p1.sub(p2.sub(p1).mul(x / y));
54 P q2 = p2.sub(p1).mul(sqrt(d) / y);
55 return new P[]{q1.sub(q2), q1.add(q2)};
56 }
57 //二円の交点
58 P[] isCC(P c1, double r1, P c2, double r2) {
59 double x = c1.sub(c2).abs2();
60 double y = ((r1 * r1 - r2 * r2) / x + 1) / 2;
61 double d = r1 * r1 / x - y * y;
62 if (d < -EPS) return new P[0];
63 if (d < 0) d = 0;
64 P q1 = c1.add(c2.sub(c1).mul(y));
65 P q2 = c2.sub(c1).mul(sqrt(d)).rot90();
66 return new P[]{q1.sub(q2), q1.add(q2)};
67 }
68 //点 p をから引いた接線の接点
69 P[] tanCP(P c, double r, P p) {
70 double x = p.sub(c).abs2();
71 double d = x - r * r;
72 if (d < -EPS) return new P[0];
73 if (d < 0) d = 0;
74 P q1 = p.sub(c).mul(r * r / x);
75 P q2 = p.sub(c).mul(-r * sqrt(d) / x).rot90();
76 return new P[]{c.add(q1.sub(q2)), c.add(q1.add(q2))};
77 }
78 //二円の共通接線
79 P[][] tanCC(P c1, double r1, P c2, double r2) {
80 List<P[]> list = new ArrayList<P[]>();
81 if (abs(r1 - r2) < EPS) {
82 P dir = c2.sub(c1);
83 dir = dir.mul(r1 / dir.abs()).rot90();
84 list.add(new P[] {c1.add(dir), c2.add(dir)});
85 list.add(new P[] {c1.sub(dir), c2.sub(dir)});
86 } else {
87 P p = c1.mul(-r2).add(c2.mul(r1)).div(r1 - r2);
88 P[] ps = tanCP(c1, r1, p);
89 P[] qs = tanCP(c2, r2, p);
90 for (int i = 0; i < ps.length && i < qs.length; i++) {
91 list.add(new P[] {ps[i], qs[i]});
92 }
93 }
94 P p = c1.mul(r2).add(c2.mul(r1)).div(r1 + r2);
95 P[] ps = tanCP(c1, r1, p);
96 P[] qs = tanCP(c2, r2, p);
97 for (int i = 0; i < ps.length && i < qs.length; i++) {
98 list.add(new P[] {ps[i], qs[i]});
99 }
100 return list.toArray(new P[0][]);
101 }
102 //２円の共通部分面積
103 double areaCC(P c1, double r1, P c2, double r2) {
104 double d = c1.sub(c2).abs();
105 if (r1 + r2 < d + EPS) return 0;
106 if (d < abs(r1 - r2) + EPS) {
107 double r = min(r1, r2);
108 return r * r * PI;
109 }
110 double x = (d * d + r1 * r1 - r2 * r2) / (2 * d);
111 double t1 = acos(x / r1);
112 double t2 = acos((d - x) / r2);
113 return r1 * r1 * t1 + r2 * r2 * t2 - d * r1 * sin(t1);
114 }
115 //原点中心半径 r の円と原点と p1,p2 からなる三角形の共通部分面積
116 double areaCT(double r, P p1, P p2) {
117 P[] qs = isCL(O, r, p1, p2);
118 if (qs.length == 0) return r * r * rad(p1, p2) / 2;
119 boolean b1 = p1.abs() > r + EPS, b2 = p2.abs() > r + EPS;
120 if (b1 && b2) {
121 if (p1.sub(qs[0]).dot(p2.sub(qs[0])) < EPS &&
122 p1.sub(qs[1]).dot(p2.sub(qs[1])) < EPS) {
123 return (r * r * (rad(p1, p2) - rad(qs[0], qs[1])) + qs[0].det(qs[1])) / 2;
124 } else {
125 return r * r * rad(p1, p2) / 2;
126 }
127 } else if (b1) {
128 return (r * r * rad(p1, qs[0]) + qs[0].det(p2)) / 2;
129 } else if (b2) {
130 return (r * r * rad(qs[1], p2) + p1.det(qs[1])) / 2;
131 } else {
132 return p1.det(p2) / 2;
133 }
134 }
\end{Code}
\subsection{4.2 Bao lồi}
\begin{NoteBlock}
Tạo bao lồi theo chiều ngược kim đồng hồ.
Không bao gồm điểm nằm trên cạnh.
Có thể chọn cách xử lý điểm trên cạnh bằng cách đổi dấu bất đẳng thức.
Nếu bao gồm điểm trên cạnh thì không được có điểm trùng.
\end{NoteBlock}
\begin{Code}
1 P[] convexHull(P[] ps) {
2 int n = ps.length, k = 0;
3 if (n <= 1) return ps;
4 sort(ps);
5 P[] qs = new P[n * 2];
6 for (int i = 0; i < n; qs[k++] = ps[i++]) {
7 while (k > 1 && qs[k - 1].sub(qs[k - 2]).det(ps[i].sub(qs[k - 1])) < EPS) k--;
8 }
9 for (int i = n - 2, t = k; i >= 0; qs[k++] = ps[i--]) {
10 while (k > t && qs[k - 1].sub(qs[k - 2]).det(ps[i].sub(qs[k - 1])) < EPS) k--;
11 }
12 P[] res = new P[k - 1];
13 System.arraycopy(qs, 0, res, 0, k - 1);
14 return res;
15 }
\end{Code}
\subsection{4.3 Cắt đa giác lồi}
\begin{NoteBlock}
O(n)
Cắt đa giác lồi bởi đường thẳng p1p2 và giữ phía bên trái.
\end{NoteBlock}
\begin{Code}
1 P[] convexCut(P[] ps, P p1, P p2) {
2 int n = ps.length;
3 ArrayList<P> res = new ArrayList<P>();
4 for (int i = 0; i < n; i++) {
5 int d1 = sig(p2.sub(p1).det(ps[i].sub(p1)));
6 int d2 = sig(p2.sub(p1).det(ps[(i + 1) % n].sub(p1)));
7 if (d1 >= 0) res.add(ps[i]);
8 if (d1 * d2 < 0) res.add(isLL(p1, p2, ps[i], ps[(i + 1) % n]));
9 }
10 return res.toArray(new P[0]);
11 }
\end{Code}
\subsection{4.4 Arrangement của các đoạn thẳng}
\begin{NoteBlock}
Lấy đầu mút và giao điểm của đoạn thẳng làm đỉnh; nếu hai điểm cùng nằm trên một đoạn và giữa chúng không có điểm khác, dùng khoảng cách giữa chúng làm trọng số
để tạo cạnh của đồ thị.
O(n3
)
\end{NoteBlock}
\begin{Code}
1 V[] segmentArrangement(P[] ps1, P[] ps2) {
2 int n = ps1.length;
3 TreeSet<V> set = new TreeSet<V>();
4 for (int i = 0; i < n; i++) {
5 set.add(new V(ps1[i]));
6 set.add(new V(ps2[i]));
7 for (int j = 0; j < i; j++) {
8 //線分交差判定のうち、端点が他の線分上の条件は抜くこと
9 if (crsSS(ps1[i], ps2[i], ps1[j], ps2[j])) {
10 set.add(new V(isLL(ps1[i], ps2[i], ps1[j], ps2[j])));
11 }
12 }
13 }
14 V[] vs = set.toArray(new V[0]);
15 for (int i = 0; i < n; i++) {
16 ArrayList<V> list = new ArrayList<V>();
17 for (V v : vs) {
18 if (crsSP(ps1[i], ps2[i], v.p)) list.add(v);
19 }
20 V[] us = list.toArray(new V[0]);
21 sort(us);
22 for (int j = 0; j + 1 < us.length; j++) {
23 us[j].connect(us[j + 1]);
24 }
25 }
26 return vs;
27 }
\end{Code}
\subsection{4.5 Kiểm tra điểm nằm trong/ngoài đa giác}
\begin{NoteBlock}
O(n)
Trả về 1 nếu ở trong, 0 nếu nằm trên cạnh, và -1 nếu ở ngoài.
\end{NoteBlock}
\begin{Code}
1 int contains(P[] ps, P q) {
2 int n = ps.length;
3 int res = -1;
4 for (int i = 0; i < n; i++) {
5 P a = ps[i].sub(q), b = ps[(i + 1) % n].sub(q);
6 if (a.y > b.y) {
7 P t = a; a = b; b = t;
8 }
9 if (a.y < EPS && b.y > EPS && a.det(b) > EPS) {
10 res = -res;
11 }
12 if (abs(a.det(b)) < EPS && a.dot(b) < EPS) return 0;
13 }
14 return res;
15 }
\end{Code}
\subsection{4.6 Khoảng cách từ đa giác lồi đến điểm bên ngoài}
\begin{NoteBlock}
Chia mặt phẳng bằng các pháp tuyến của cạnh đa giác lồi, rồi dùng tìm kiếm nhị phân để xác định miền chứa điểm.
O(log n)
Giả sử đa giác lồi được cho theo chiều ngược kim đồng hồ.
\end{NoteBlock}
\begin{Code}
1 double disConvexP(P[] ps, P q) {
2 int n = ps.length;
3 int left = 0, right = n;
4 while (right - left > 1) {
5 int mid = (left + right) / 2;
6 if (in(ps[(left + n - 1) % n], ps[left], ps[mid], ps[(mid + 1) % n], q)) {
7 right = mid;
8 } else {
9 left = mid;
10 }
11 }
12 return disSP(ps[left], ps[right % n], q);
13 }
14 boolean in(P p1, P p2, P p3, P p4, P q) {
15 P o12 = p1.sub(p2).rot90();
16 P o23 = p2.sub(p3).rot90();
17 P o34 = p3.sub(p4).rot90();
18 return in(o12, o23, q.sub(p2)) || in(o23, o34, q.sub(p3))
19 || in(o23, p3.sub(p2), q.sub(p2)) && in(p2.sub(p3), o23, q.sub(p3));
20 }
21 boolean in(P p1, P p2, P q) {
22 return p1.det(q) > -EPS && p2.det(q) < EPS;
23 }
\end{Code}
\subsection{4.7 Đường kính của đa giác lồi}
\begin{NoteBlock}
Dùng rotating calipers để tìm khoảng cách xa nhất giữa hai điểm trên đa giác lồi.
O(n)
\end{NoteBlock}
\begin{Code}
1 double convexDiameter(P[] ps) {
2 int n = ps.length;
3 int is = 0, js = 0;
4 for (int i = 1; i < n; i++) {
5 if (ps[i].x > ps[is].x) is = i;
6 if (ps[i].x < ps[js].x) js = i;
7 }
8 double maxd = ps[is].sub(ps[js]).abs();
9 int i = is, j = js;
10 do {
11 if (ps[(i + 1) % n].sub(ps[i]).det(ps[(j + 1) % n].sub(ps[j])) >= 0) {
12 j = (j + 1) % n;
13 } else {
14 i = (i + 1) % n;
15 }
16 maxd = max(maxd, ps[i].sub(ps[j]).abs());
17 } while (i != is || j != js);
18 return maxd;
19 }
\end{Code}
\subsection{4.8 Tập công thức}
\subsubsection{4.8.1 Tâm nội tiếp tam giác}
\begin{NoteBlock}
a
-
⇀
A + b
-
⇀
B + c
-
⇀
C
a + b + c
\end{NoteBlock}
\subsubsection{4.8.2 Tâm ngoại tiếp tam giác}
\begin{NoteBlock}
-
⇀
A +
-
⇀
B -
-
-
⇀
BC·
-
⇀
CA
-
-
⇀
AB×
-
-
⇀
BC
-
-
⇀
ABT
2
\end{NoteBlock}
\subsubsection{4.8.3 Trực tâm tam giác}
\begin{NoteBlock}
-
⇀
H = 3
-
⇀
G - 2
-
⇀
O
\end{NoteBlock}
\subsubsection{4.8.4 Tâm bàng tiếp tam giác}
\begin{NoteBlock}
-a
-
⇀
A + b
-
⇀
B + c
-
⇀
C
-a + b + c
Hai điểm còn lại cũng tương tự.
\end{NoteBlock}
\subsubsection{4.8.5 Bán kính đường tròn ngoại tiếp tam giác}
\begin{NoteBlock}
R =
abc
4S
\end{NoteBlock}
\subsubsection{4.8.6 Công thức Heron}
\begin{NoteBlock}
2s := a + b + c
S =
√
s(s - a)(s - b)(s - c)
\end{NoteBlock}
\subsubsection{4.8.7 Công thức Pick}
\begin{NoteBlock}
Với đa giác đơn có mọi đỉnh nằm trên điểm nguyên, nếu S là diện tích, B là số điểm nguyên trên biên, I là số điểm nguyên bên trong, thì:
S =
B
2
+ I - 1
\end{NoteBlock}
\section{5 Hình học không gian}
\subsection{5.1 Giao điểm và các phép tính liên quan}
\begin{Code}
1 //直線と点の距離
2 double disLP(P p1, P p2, P q) {
3 return p2.sub(p1).det(q.sub(p1)).abs() / p2.sub(p1).abs();
4 }
5 //直線と直線の距離
6 double disLL(P p1, P p2, P q1, P q2) {
7 P p = q1.sub(p1);
8 P u = p2.sub(p1);
9 P v = q2.sub(q1);
10 double d = u.abs2() * v.abs2() - u.dot(v) * u.dot(v);
11 if (abs(d) < EPS) return disLP(q1, q2, p1);
12 double s = (p.dot(u) * v.abs2() - p.dot(v) * u.dot(v)) / d;
13 return disLP(q1, q2, p1.add(u.mul(s)));
14 }
15 //平面と直線の交点
16 P isFL(P p, P o, P q1, P q2) {
17 double a = o.dot(q2.sub(p));
18 double b = o.dot(q1.sub(p));
19 double d = a - b;
20 if (abs(d) < EPS) return null;
21 return q1.mul(a).sub(q2.mul(b)).div(d);
22 }
23 //平面と平面の交線
24 P[] isFF(P p1, P o1, P p2, P o2) {
25 P e = o1.det(o2);
26 P v = o1.det(e);
27 double d = o2.dot(v);
28 if (abs(d) < EPS) return null;
29 P q = p1.add(v.mul(o2.dot(p2.sub(p1)) / d));
30 return new P[]{q, q.add(e)};
31 }
\end{Code}
\subsection{5.2 Bao lồi không gian}
\begin{NoteBlock}
Trả về các mặt sao cho nhìn từ bên ngoài thì theo chiều ngược kim đồng hồ.
Cần nhiễu để không có bốn điểm đồng phẳng.
O(n2
)
\end{NoteBlock}
\begin{Code}
1 P[][] convexHull(P[] ps) {
2 int n = ps.length;
3 //vs[i][j]=辺 i->j をこの向きに含む三角形が、まだ調べてない:0、残った:-1、取り除かれた:1
4 int[][] vs = new int[n][n];
5 List<int[]> crt = new ArrayList<int[]>();
6 crt.add(new int[]{0, 1, 2});
7 crt.add(new int[]{2, 1, 0});
8 for (int i = 3; i < n; i++) {
9 List<int[]> next = new ArrayList<int[]>();
10 for (int[] t : crt) {
11 int v = ps[t[1]].sub(ps[t[0]]).det(ps[t[2]].sub(ps[t[0]])).dot(
12 ps[i].sub(ps[t[0]])) < 0 ? -1 : 1;
13 if (v < 0) next.add(t);
14 for (int j = 0; j < 3; j++) {
15 if (vs[t[(j + 1) % 3]][t[j]] == 0) {
16 vs[t[j]][t[(j + 1) % 3]] = v;
17 } else {
18 if (vs[t[(j + 1) % 3]][t[j]] != v) {
19 if (v > 0) next.add(new int[]{t[j], t[(j + 1) % 3], i});
20 else next.add(new int[]{t[(j + 1) % 3], t[j], i});
21 }
22 vs[t[(j + 1) % 3]][t[j]] = 0;
23 }
24 }
25 }
26 crt = next;
27 }
28 P[][] pss = new P[crt.size()][3];
29 for (int i = 0; i < pss.length; i++) {
30 for (int j = 0; j < 3; j++) pss[i][j] = ps[crt.get(i)[j]];
31 }
32 return pss;
33 }
\end{Code}
\subsection{5.3 Cắt đa diện lồi}
\begin{NoteBlock}
O(n log n)
Cắt đa diện lồi bằng mặt phẳng po và giữ phía sau.
Đa diện lồi phải được định hướng ngược kim đồng hồ khi nhìn từ ngoài.
\end{NoteBlock}
\begin{Code}
1 P[][] convexCut(P[][] pss, P p, P o) {
2 ArrayList<P[]> res = new ArrayList<P[]>();
3 ArrayList<P> sec = new ArrayList<P>();
4 for (P[] ps : pss) {
5 int n = ps.length;
6 ArrayList<P> qs = new ArrayList<P>();
7 boolean dif = false;
8 for (int i = 0; i < n; i++) {
9 int d1 = sig(o.dot(ps[i].sub(p)));
10 int d2 = sig(o.dot(ps[(i + 1) % n].sub(p)));
11 if (d1 <= 0) qs.add(ps[i]);
12 if (d1 * d2 < 0) {
13 P q = isFL(p, o, ps[i], ps[(i + 1) % n]);
14 qs.add(q);
15 sec.add(q);
16 }
17 if (d1 == 0) sec.add(ps[i]);
18 else dif = true;
19 dif |= o.dot(ps[(i + 1) % n].sub(ps[i]).det(ps[(i + 2) % n].sub(ps[i]))) < -EPS;
20 }
21 if (qs.size() > 0 && dif) res.add(qs.toArray(new P[0]));
22 }
23 //2D の convexHull に.dot(o) をつけたもの
24 if (sec.size() > 0) res.add(convexHull2D(sec.toArray(new P[0]), o));
25 return res.toArray(new P[0][]);
26 }
27 //空間の初期化
28 P[][] init() {
29 P[][] pss = new P[6][4];
30 pss[0][0] = pss[1][0] = pss[2][0] = new P(-INF, -INF, -INF);
31 pss[0][3] = pss[1][1] = pss[5][2] = new P(-INF, -INF, INF);
32 pss[0][1] = pss[2][3] = pss[4][2] = new P(-INF, INF, -INF);
33 pss[0][2] = pss[5][3] = pss[4][1] = new P(-INF, INF, INF);
34 pss[1][3] = pss[2][1] = pss[3][2] = new P( INF, -INF, -INF);
35 pss[1][2] = pss[5][1] = pss[3][3] = new P( INF, -INF, INF);
36 pss[2][2] = pss[4][3] = pss[3][1] = new P( INF, INF, -INF);
37 pss[5][0] = pss[4][0] = pss[3][0] = new P( INF, INF, INF);
38 return pss;
39 }
\end{Code}
\subsection{5.4 Tập công thức}
\subsubsection{5.4.1 Công thức tứ diện Euler}
\begin{NoteBlock}
Thể tích tứ diện có sáu cạnh dài a, b, c, d, e, f.
Ghi chú/công thức: O - ABC   a = AB, b = BC, c = CA, d = OC, e = OA, f = OB
(12V )2
= a2
d2
(b2
+ c2
+ e2
+ f2
- a2
- d2
) + b2
e2
(c2
+ a2
+ f2
+ d2
- b2
- e2
) + c2
f2
(a2
+ b2
+ d2
+ e2
-
c2
- f2
) - a2
b2
c2
- a2
e2
f2
- d2
b2
f2
- d2
e2
c2
\end{NoteBlock}
\subsection{5.5 Phép biến đổi}
\subsubsection{5.5.1 Biến đổi đối ngẫu}
\begin{NoteBlock}
Phương pháp biến đổi.
Ghi chú/công thức: d   p = (p1, p2, . . . , pd)   d   p∗
Ghi chú/công thức: : xd = p1x1 + p2x2 + . . . + pd-1xd-1 - pd
Ghi chú/công thức: 
Ghi chú/công thức: d   h : xd = a1x1 + a2x2 + . . . + ad   d   h∗
Ghi chú/công thức: = (a1, a2, . . . , ad-1, -ad)
Tính chất.
Ghi chú/công thức: 
Nói cách khác:
p ∈ h ⇔ h∗
∈ p∗
Ghi chú/công thức: p   h   ⇔ h∗
Ghi chú/công thức: p∗
Ghi chú/công thức: 
Envelope và bao lồi.
Ghi chú/công thức: h   H   h∗
Ghi chú/công thức: H∗
Ghi chú/công thức: 
Ghi chú/công thức: 
\end{NoteBlock}
\subsubsection{5.5.2 Biến đổi nâng/inversion 1}
\begin{NoteBlock}
Phương pháp biến đổi.
Ghi chú/công thức: p = (x, y)   p′
= (x, y, x2
+ y2
Ghi chú/công thức: )
Ghi chú/công thức: C : (x - a)2
+ (y - b)2
= r2
Ghi chú/công thức: C′
: z = a(x - a) + b(y - b) + r2
Ghi chú/công thức: 
Tính chất.
Ghi chú/công thức: 
Nói cách khác:
Ghi chú/công thức: p   C   ⇔ p′
Ghi chú/công thức: C′
Ghi chú/công thức: 
Phân hoạch Delaunay và biểu đồ Voronoi.
Ghi chú/công thức: P   P′
Ghi chú/công thức: 
P′
Ghi chú/công thức: z = 0   P  P′∗
Ghi chú/công thức: 
Ghi chú/công thức: z = 0   P
\end{NoteBlock}
\subsubsection{5.5.3 Biến đổi inversion 2}
\begin{NoteBlock}
Phương pháp biến đổi.
Ghi chú/công thức: p = (r, θ)   p′
= (1
Ghi chú/công thức: r , θ)
Ghi chú/công thức: 
Tính chất.
Ghi chú/công thức: 
\end{NoteBlock}
\section{6 Toán học}
\subsection{6.1 Đếm tổ hợp}
\subsubsection{6.1.1 Số Catalan}
\begin{NoteBlock}
C(n) =
2nCn
n + 1
Số cách tạo cây nhị phân với n nút là C(n).
Ghi chú/công thức: n   n - 2   C(n - 2)
\end{NoteBlock}
\subsubsection{6.1.2 Tổ hợp có lặp}
\begin{NoteBlock}
Tổng số cách chọn k phần tử từ n loại, trong đó loại i có m[i] phần tử.
x[i, j] = x[i - 1, j] + x[i, j - 1] - x[i - 1, j - m[i] - 1]
\end{NoteBlock}
\subsubsection{6.1.3 Số phân hoạch}
\begin{NoteBlock}
Số cách phân hoạch n thành k số nguyên không âm.
Ghi chú/công thức: P(n - k, k)  P(n, n)
P[i, j] = P[i - 1, j] + P[i, j - i]
Phân hoạch thành số lượng tùy ý các số tự nhiên có thể tính nhanh hơn.
O(n
3
\end{NoteBlock}
\begin{Code}
2 )
1 int partition(int n) {
2 int[] dp = new int[n + 1];
3 dp[0] = 1;
4 for (int i = 1; i <= n; i++) {
5 for (int j = 1, r = 1; i - (3 * j * j - j) / 2 >= 0; j++, r *= -1) {
6 dp[i] += dp[i - (3 * j * j - j) / 2] * r;
7 if (i - (3 * j * j + j) / 2 >= 0) {
8 dp[i] += dp[i - (3 * j * j + j) / 2] * r;
9 }
10 }
11 }
12 return dp[n];
13 }
\end{Code}
\subsubsection{6.1.4 Số Bell}
\begin{NoteBlock}
Tổng số cách chia n phần tử phân biệt thành các nhóm.
B[n] =
∑n-1
i=0 C[n - 1, i]B[i]
\end{NoteBlock}
\subsubsection{6.1.5 Số Stirling loại một}
\begin{NoteBlock}
Số cách chia n phần tử phân biệt thành k chu trình không rỗng.
Đồng thời là hệ số khi biểu diễn lũy thừa giai thừa tăng bằng tổng các lũy thừa thường.
S[n, k] = (n - 1)S[n - 1, k] + S[n - 1, k - 1]
\end{NoteBlock}
\subsubsection{6.1.6 Số Stirling loại hai}
\begin{NoteBlock}
Số cách chia n phần tử phân biệt thành k nhóm không rỗng.
Đồng thời là hệ số khi biểu diễn lũy thừa thường bằng tổng các lũy thừa giai thừa giảm.
B[n] =
∑n
Ghi chú/công thức: i=1 S[n, i]
S[n, k] = kS[n - 1, k] + S[n - 1, k - 1]
Khi chia thành các nhóm có ít nhất r phần tử:
S[n, k] = kS[n - 1, k] + C[n - 1, r - 1]S[n - r, k - 1]
ta có công thức sau.
\end{NoteBlock}
\subsubsection{6.1.7 Số Euler}
\begin{NoteBlock}
Số hoán vị trên {1,...,n} có đúng k lần tăng.
E[n, k] = (k + 1)E[n - 1, k] + (n - k)E[n - 1, k - 1]
\end{NoteBlock}
\subsubsection{6.1.8 Nguyên lý bao hàm - loại trừ}
\begin{NoteBlock}
∪
P ∈A
P =
∑
∅⊂S⊆A
(-1)|S|-1
∩
P ∈S
P
g(A) =
∑
S⊆A
f(S) ⇐⇒ f(A) =
∑
S⊆A
(-1)|A|-|S|
g(S)
\end{NoteBlock}
\subsubsection{6.1.9 Hàm Mobius}
\begin{NoteBlock}
µ(n) =
{
0 (khi n có thừa số chính phương)
(-1)k
Ghi chú/công thức: (n   k  )
g(n) =
∑
d|n
f(d) ⇐⇒ f(n) =
∑
d|n
µ(d)g(
n
d
)
g(x) =
⌊x⌋
∑
n=1
f(
x
n
) ⇐⇒ f(x) =
⌊x⌋
∑
n=1
µ(n)g(
x
n
)
\end{NoteBlock}
\subsubsection{6.1.10 Bổ đề Burnside, tên cũ}
\begin{NoteBlock}
|X/G| =
1
|G|
∑
g∈G
|Xg|
\end{NoteBlock}
\subsection{6.2 Lý thuyết số}
\subsubsection{6.2.1 Tập công thức}
\begin{NoteBlock}
Định lý nhỏ Fermat.
ap
≡ a (mod p)
Ghi chú/công thức: (a, p) = 1   ap-1
≡ 1 (mod p)
Ghi chú/công thức: a-1
≡ ap-2
(mod p)
Hàm Euler phi.
phi(n) là số số tự nhiên từ 1 đến n nguyên tố cùng nhau với n.
ϕ(n) = n
∏
p|n
(1 -
1
p
)
Ghi chú/công thức: (a, n) = 1   aϕ(n)
≡ 1 (mod n)
Bậc.
ax
Ghi chú/công thức: ≡ 1 (mod m)   x   a (mod m)
Ghi chú/công thức: a (mod m)  ϕ(m)
a được gọi là căn nguyên thủy modulo m.
Ghi chú/công thức: p  ϕ(p - 1)
xd
Ghi chú/công thức: ≡ 1 (mod p)
Ghi chú/công thức: d | (p - 1)   (d  ) ϕ(d)   d
Định lý Wilson.
(p - 1)! ≡ -1 (mod p)
Thặng dư bậc hai.
Với số nguyên tố lẻ p:
x2
Ghi chú/công thức: ≡ a (mod p)   ⇐⇒ a
p-1
2 ≡ 1 (mod p)
Cây Stern-Brocot.
Đây là cây tìm kiếm nhị phân có các khoảng hữu tỉ làm đỉnh, được định nghĩa như sau.
Gốc là:
(
0
1
,
1
0
)
Với đỉnh
(a
b
,
c
d
)
các con là
(
a
b
,
a + c
b + d
)
Ghi chú/công thức: 
(
a + c
b + d
,
c
d
)
Cây này có các tính chất sau.
Các đầu mút của mỗi đỉnh là phân số tối giản.
Mọi phân số tối giản xuất hiện đúng một lần.
Theo độ sâu của cây, tử số và mẫu số tăng đơn điệu.
Với đỉnh
(a
b
,
c
d
)
với đỉnh đó,
ad - bc = 1
Công thức Stirling.
n! ≈
√
2πn
nn
en
■nCk mod p
Ghi chú/công thức: nCk  p   ⇐⇒ k + (n - k)   p
Ghi chú/công thức: n =
∑
nipi
, k =
∑
kipi
Ghi chú/công thức: 
nCk ≡
∏
ni Cki (mod p)
\end{NoteBlock}
\subsubsection{6.2.2 Giai thừa modulo}
\begin{NoteBlock}
Với số nguyên tố p, viết n! = a p^k.
Tìm giá trị a modulo p.
O(p log n)
\end{NoteBlock}
\begin{Code}
1 int modFact(int n, int p) {
2 int res = 1;
3 while (n > 0) {
4 for (int i = 1, m = n % p; i <= m; i++) res = (res * i) % p;
5 if ((n /= p) % 2 > 0) res = p - res;
6 }
7 return res;
8 }
\end{Code}
\subsubsection{6.2.3 Thuật toán Euclid mở rộng}
\begin{NoteBlock}
Tìm một bộ {a,b,c=gcd(x,y)} sao cho ax + by = gcd(x,y).
Ghi chú/công thức: (a, b)   (a + dy
c , b - dx
c )
\end{NoteBlock}
\begin{Code}
1 int[] exgcd(int x, int y) {
2 int a0 = 1, a1 = 0, b0 = 0, b1 = 1, t;
3 while (y != 0) {
4 t = a0 - x / y * a1; a0 = a1; a1 = t;
5 t = b0 - x / y * b1; b0 = b1; b1 = t;
6 t = x % y; x = y; y = t;
7 }
8 if (x < 0) {
9 a0 = -a0; b0 = -b0; x = -x;
10 }
11 return new int[] {a0, b0, x};
12 }
\end{Code}
\subsubsection{6.2.4 Hệ đồng dư tuyến tính}
\begin{NoteBlock}
Trả về nghiệm và modulo {x,m} của hệ đồng dư tuyến tính Ax == B (mod M).
Nếu không tồn tại nghiệm thì trả về null.
\end{NoteBlock}
\begin{Code}
1 BigInt[] congruence(BigInt[] A, BigInt[] B, BigInt[] M) {
2 BigInt x = 0, m = 1;
3 for (int i = 0; i < A.length; i++) {
4 BigInt a = A[i] * m, b = B[i] - A[i] * x, d = a.gcd(M[i]);
5 if (b % d != 0) return null;
6 x += m * (b / d * (a / d).modInv(M[i] / d) % (M[i] / d));
7 m *= M[i] / d;
8 }
9 return new BigInt[] {x % m, m};
10 }
\end{Code}
\subsubsection{6.2.5 Logarit rời rạc}
\begin{NoteBlock}
ax
Ghi chú/công thức: ≡ b (mod m)   x
Ghi chú/công thức: -1
\end{NoteBlock}
\begin{Code}
1 long modLog(long a, long b, long m) {
2 if (b % exgcd(a, m)[2] != 0) return -1;
3 if (m == 0) return 0;
4 long n = (long)sqrt(m) + 1;
5 Map<Long, Long> map = new HashMap<Long, Long>();
6 long an = 1;
7 for (long j = 0; j < n; j++) {
8 if (!map.containsKey(an)) map.put(an, j);
9 an = an * a % m;
10 }
11 long ain = 1, res = Long.MAX_VALUE;
12 for (long i = 0; i < n; i++) {
13 long[] is = congruence(ain, b, m);
14 for (long aj = is[0]; aj < m; aj += is[1]) if (map.containsKey(aj)) {
15 long j = map.get(aj);
16 res = min(res, i * n + j);
17 }
18 if (res < Long.MAX_VALUE) return res;
19 ain = ain * an % m;
20 }
21 return -1;
22 }
\end{Code}
\subsubsection{6.2.6 Giải đồng dư modulo hợp số}
\begin{NoteBlock}
Dạng tuyến tính.
ax + by = n có nghiệm khi và chỉ khi gcd(a,b) chia n.
ka ≡ kb (mod m) ⇐⇒ a ≡ b (mod
m
(k, m)
)
kx == l (mod m) có gcd(k,m) nghiệm khi gcd(k,m) chia l.
Đa thức.
Muốn tìm nghiệm của f(x) == 0 (mod m).
m =
k
∏
i=1
pai
Ghi chú/công thức: i
Ghi chú/công thức: f(x) ≡ 0 (mod pai
Ghi chú/công thức: i )   {xi}k
Ghi chú/công thức: i=1
Một nghiệm được xác định.
f(x) ≡ 0 (mod pa
Ghi chú/công thức: )
f(x) ≡ 0 (mod pa-1
Ghi chú/công thức: )   x′
Ghi chú/công thức: 
f′
(x′
Ghi chú/công thức: ) ≡ 0 (mod p)   f(x′
) ≡ 0 (mod pa
) ⇒ x′
+ dpa-1
(d = 0, . . . , p - 1)
f′
(x′
) ̸≡ 0 (mod p) ⇒ x′
-
f(x′
)
f′(x′)
\end{NoteBlock}
\subsection{6.3 Đại số tuyến tính}
\subsubsection{6.3.1 Tập công thức}
\begin{NoteBlock}
Công thức Cramer.
Với Ax=b, gọi A_i là ma trận thay cột i của A bằng b.
xi =
|Ai|
|A|
Hệ tuyến tính ba ẩn nên giải bằng cách này.
Biến đổi đối ngẫu.
max{cx | Ax ≤ b, x ≥ 0} = min{yb | yA ≥ c, y ≥ 0}
max{cx | Ax = b, x ≥ 0} = min{yb | yA ≥ c}
max{cx | Ax = b} = min{yb | yA = c}
\end{NoteBlock}
\subsubsection{6.3.2 Hệ phương trình}
\begin{NoteBlock}
Ghi chú/công thức: ■Ax = 0   X = {x0 +
∑
Ghi chú/công thức: aixi}   O(nmr)
\end{NoteBlock}
\begin{Code}
1 double[][] solutionSpace(double[][] A, double[] b) {
2 int n = A.length, m = A[0].length;
3 double[][] a = new double[n][m + 1];
4 for (int i = 0; i < n; i++) {
5 System.arraycopy(A[i], 0, a[i], 0, m);
6 a[i][m] = b[i];
7 }
8 int[] id = new int[n + 1];//単位行列の成分が (i,id[i]) になる
9 fill(id, -1);
10 int pi = 0;//ランクに相当
11 for (int pj = 0; pi < n && pj < m; pj++) {
12 for (int i = pi + 1; i < n; i++) {
13 if (abs(a[i][pj]) > abs(a[pi][pj])) {
14 double[] t = a[i]; a[i] = a[pi]; a[pi] = t;
15 }
16 }
17 if (abs(a[pi][pj]) < EPS) continue;
18 double inv = 1 / a[pi][pj];
19 for (int j = 0; j <= m; j++) a[pi][j] *= inv;
20 for (int i = 0; i < n; i++) if (i != pi) {
21 double d = a[i][pj];
22 for (int j = 0; j <= m; j++) a[i][j] -= d * a[pi][j];
23 }
24 id[pi++] = pj;
25 }
26 for (int i = pi; i < n; i++) if (abs(a[i][m]) > EPS) return null;
27 double[][] X = new double[1 + m - pi][m];
28 for (int j = 0, k = 0; j < m; j++) {
29 if (id[k] == j) X[0][j] = a[k++][m];
30 else {
31 for (int i = 0; i < k; i++) X[1 + j - k][id[i]] = -a[i][j];
32 X[1 + j - k][j] = 1;
33 }
34 }
35 return X;
36 }
\end{Code}
\begin{NoteBlock}
Ghi chú/công thức: ■Toeplitz   A  Ax = b   O(n2
)
A[i, j] = A[i - j + ZERO]
A[0, 0] ̸= 0
\end{NoteBlock}
\begin{Code}
1 double[] levinson(double[] A, double[] b) {
2 int n = (A.length + 1) / 2, ZERO = n - 1;
3 double[] x = new double[n], fs = new double[n], bs = new double[n];
4 fs[0] = bs[0] = 1 / A[ZERO];
5 x[0] = b[0] / A[ZERO];
6 for (int i = 1; i < n; i++) {
7 double ef = 0, eb = 0, ex = 0;
8 for (int j = 0; j < i; j++) {
9 ef += fs[j] * A[i - j + ZERO];
10 eb += bs[j] * A[-1 - j + ZERO];
11 ex += x[j] * A[i - j + ZERO];
12 }
13 if (abs(ef * eb - 1) < EPS) return null;
14 for (int j = i; j >= 0; j--) {
15 double af = (j < i ? fs[j] : 0), ab = (j != 0 ? bs[j - 1] : 0);
16 fs[j] = af / (1 - ef * eb) - ef * ab / (1 - ef * eb);
17 bs[j] = ab / (1 - ef * eb) - eb * af / (1 - ef * eb);
18 }
19 for (int j = 0; j <= i; j++) x[j] += (b[i] - ex) * bs[j];
20 }
21 return x;
22 }
\end{Code}
\subsubsection{6.3.3 Phương pháp đơn hình}
\begin{NoteBlock}
Giải max{cx | Ax <= b, x >= 0}.
Nếu không tồn tại nghiệm thì trả về null.
\end{NoteBlock}
\begin{Code}
1 double[] simplex(double[][] A, double[] b, double[] c) {
2 int n = A.length, m = A[0].length + 1, r = n, s = m - 1;
3 double[][] D = new double[n + 2][m + 1];
4 int[] ix = new int[n + m];
5 for (int i = 0; i < n + m; i++) ix[i] = i;
6 for (int i = 0; i < n; i++) {
7 for (int j = 0; j < m - 1; j++) D[i][j] = -A[i][j];
8 D[i][m - 1] = 1;
9 D[i][m] = b[i];
10 if (D[r][m] > D[i][m]) r = i;
11 }
12 for (int j = 0; j < m - 1; j++) D[n][j] = c[j];
13 D[n + 1][m - 1] = -1;
14 for (double d;;) {
15 if (r < n) {
16 int t = ix[s]; ix[s] = ix[r + m]; ix[r + m] = t;
17 D[r][s] = 1.0 / D[r][s];
18 for (int j = 0; j <= m; j++) if (j != s) D[r][j] *= -D[r][s];
19 for (int i = 0; i <= n + 1; i++) if (i != r) {
20 for (int j = 0; j <= m; j++) if (j != s) D[i][j] += D[r][j] * D[i][s];
21 D[i][s] *= D[r][s];
22 }
23 }
24 r = -1; s = -1;
25 for (int j = 0; j < m; j++) if (s < 0 || ix[s] > ix[j]) {
26 if (D[n + 1][j] > EPS || D[n + 1][j] > -EPS && D[n][j] > EPS) s = j;
27 }
28 if (s < 0) break;
29 for (int i = 0; i < n; i++) if (D[i][s] < -EPS) {
30 if (r < 0 || (d = D[r][m] / D[r][s] - D[i][m] / D[i][s]) < -EPS
31 || d < EPS && ix[r + m] > ix[i + m]) r = i;
32 }
33 if (r < 0) return null;//非有界
34 }
35 if (D[n + 1][m] < -EPS) return null;//実行不能
36 double[] x = new double[m - 1];
37 for (int i = m; i < n + m; i++) if (ix[i] < m - 1) x[ix[i]] = D[i - m][m];
38 return x;//値は D[n][m]
39 }
\end{Code}
\subsubsection{6.3.4 Công thức truy hồi}
\begin{NoteBlock}
Ghi chú/công thức: ai+n =
∑
Ghi chú/công thức: kjai+j + d  am =
∑
Ghi chú/công thức: ciai + cnd   {ci}
O(n2
log m)
Ghi chú/công thức: {ki}   FFT   O(n log n log m)
\end{NoteBlock}
\begin{Code}
1 double[] recFormula(double[] k, long m) {
2 int n = k.length;
3 double[] c = new double[n + 1];
4 if (m < n) c[(int)m] = 1;
5 else {
6 double[] b = recFormula(k, m >>> 1);
7 double[] a = new double[n * 2];
8 int s = (int)(m & 1);
9 for (int i = 0; i < n; i++) {
10 for (int j = 0; j < n; j++) {
11 a[i + j + s] += b[i] * b[j];
12 }
13 c[n] += b[i];
14 }
15 c[n] = (c[n] + 1) * b[n];
16 for (int i = n * 2 - 1; i >= n; i--) {
17 for (int j = 0; j < n; j++) {
18 a[i - n + j] += k[j] * a[i];
19 }
20 c[n] += a[i];
21 }
22 System.arraycopy(a, 0, c, 0, n);
23 }
24 return c;
25 }
\end{Code}
\subsubsection{6.3.5 Nghiệm nguyên}
\begin{NoteBlock}
O(n5
Ghi chú/công thức: )
Nếu không có nghiệm trả về null; nếu có nhiều nghiệm, trả về một nghiệm.
\end{NoteBlock}
\begin{Code}
1 BigInteger[] intSolve(BigInteger[][] A, BigInteger[] b) {
2 int n = A.length, m = A[0].length;
3 BigInteger[][] a = new BigInteger[n][m + 1];
4 for (int i = 0; i < n; i++) {
5 for (int j = 0; j < m; j++) a[i][j] = A[i][j];
6 a[i][m] = b[i];
7 }
8 Stack<Trans> trans = new Stack<Trans>();
9 for (int p = 0; p < m && p < n; p++) {
10 for (;;) {
11 int pi = p, pj = p;
12 for (int i = p; i < n; i++) {
13 for (int j = p; j < m; j++) {
14 if (a[i][j].signum() != 0 && (a[pi][pj].signum() == 0
15 || a[pi][pj].abs().compareTo(a[i][j].abs()) > 0)) {
16 pi = i;
17 pj = j;
18 }
19 }
20 }
21 swap(a, pi, p);
22 for (int i = p; i < n; i++) swap(a[i], pj, p);
23 if (pj != p) trans.push(new Trans(0, pj, p, null));
24 if (a[p][p].signum() == 0) break;
25 boolean end = true;
26 for (int i = p + 1; i < n; i++) {
27 BigInteger d = a[i][p].divide(a[p][p]);
28 end = end && a[i][p].signum() == 0;
29 for (int j = p; j <= m; j++) {
30 a[i][j] = a[i][j].subtract(a[p][j].multiply(d));
31 }
32 }
33 for (int j = p + 1; j < m; j++) {
34 BigInteger d = a[p][j].divide(a[p][p]);
35 end = end && a[p][j].signum() == 0;
36 trans.push(new Trans(1, j, p, d));
37 for (int i = p; i < n; i++) {
38 a[i][j] = a[i][j].subtract(a[i][p].multiply(d));
39 }
40 }
41 if (end) break;
42 }
43 }
44 BigInteger[] res = new BigInteger[m];
45 fill(res, ZERO);
46 for (int i = 0; i < m && i < n; i++) {
47 if (a[i][i].signum() < 0) {
48 a[i][i] = a[i][i].negate();
49 a[i][m] = a[i][m].negate();
50 }
51 if (a[i][i].signum() == 0) {
52 if (a[i][m].signum() != 0) return null;
53 } else if (a[i][m].mod(a[i][i]).signum() == 0) {
54 res[i] = a[i][m].divide(a[i][i]);
55 } else return null;
56 }
57 for (int i = min(m, n); i < n; i++) if (a[i][m].signum() != 0) return null;
58 while (!trans.isEmpty()) {
59 Trans t = trans.pop();
60 if (t.type == 0) swap(res, t.a, t.b);
61 else res[t.b] = res[t.b].subtract(res[t.a].multiply(t.c));
62 }
63 return res;
64 }
\end{Code}
\subsection{6.4 Giải tích}
\subsubsection{6.4.1 Biến đổi Fourier nhanh}
\begin{NoteBlock}
■ O(n log n)
Độ dài phải là lũy thừa của 2.
Dùng sign=1 cho biến đổi thuận, sign=-1 cho biến đổi ngược.
Khi biến đổi ngược, nhân kết quả với 1/n.
\end{NoteBlock}
\begin{Code}
1 void fft(int sign, double[] real, double[] imag) {
2 int n = real.length, d = Integer.numberOfLeadingZeros(n) + 1;
3 double theta = sign * 2 * PI / n;
4 for (int m = n; m >= 2; m >>= 1, theta *= 2) {
5 for (int i = 0, mh = m >> 1; i < mh; i++) {
6 double wr = cos(i * theta), wi = sin(i * theta);
7 for (int j = i; j < n; j += m) {
8 int k = j + mh;
9 double xr = real[j] - real[k], xi = imag[j] - imag[k];
10 real[j] += real[k];
11 imag[j] += imag[k];
12 real[k] = wr * xr - wi * xi;
13 imag[k] = wr * xi + wi * xr;
14 }
15 }
16 }
17 for (int i = 0; i < n; i++) {
18 int j = Integer.reverse(i) >>> d;
19 if (j < i) {
20 double tr = real[i]; real[i] = real[j]; real[j] = tr;
21 double ti = imag[i]; imag[i] = imag[j]; imag[j] = ti;
22 }
23 }
24 }
\end{Code}
\subsubsection{6.4.2 Khai triển Taylor}
\begin{NoteBlock}
ex
=
∞
∑
n=0
xn
n!
log(1 + x) =
∞
∑
n=1
(-1)n+1
n
xn
(|x| < 1)
sin x =
∞
∑
n=0
(-1)n
(2n + 1)!
x2n+1
cos x =
∞
∑
n=0
(-1)n
(2n)!
x2n
arcsin x =
∞
∑
n=0
(2n)!
4n(n!)2(2n + 1)
x2n+1
(|x| < 1)
\end{NoteBlock}
\subsubsection{6.4.3 Tích phân số}
\begin{NoteBlock}
Công thức Simpson.
Tích phân số bằng cách xấp xỉ f(x) bởi hàm bậc hai.
Integral_a^b f(x) dx ~= h/3 * [ f(x0)
  + 2 * sum_{j=1}^{n/2-1} f(x_{2j})
  + 4 * sum_{j=1}^{n/2} f(x_{2j-1})
  + f(x_n) ].
Ở đây n là số chẵn, h=(b-a)/n và x_i=a+ih.
Ghi chú/công thức: h   O(h4
)
\end{NoteBlock}
\section{7 Chuỗi}
\subsection{7.1 KMP}
\begin{NoteBlock}
Tìm mẫu p trong văn bản t.
fail[i] là độ dài lớn nhất L trong [0,i) sao cho prefix và suffix trùng nhau.
Ngoài ra, n-fail[n] biểu diễn chu kỳ của chuỗi.
Ghi chú/công thức: s  rev(s)   s  i=n-1   fail
Ghi chú/công thức: i-j
\end{NoteBlock}
\begin{Code}
1 class KMP {
2 int m;
3 char[] p;
4 int[] fail;
5 KMP(char[] p) {
6 m = p.length;
7 this.p = p;
8 fail = new int[m + 1];
9 int crt = fail[0] = -1;
10 for (int i = 1; i <= m; i++) {
11 while (crt >= 0 && p[crt] != p[i - 1]) crt = fail[crt];
12 fail[i] = ++crt;
13 }
14 }
15 int searchFrom(char[] t) {
16 int n = t.length, count = 0;
17 for (int i = 0, j = 0; i < n; i++) {
18 while (j >= 0 && t[i] != p[j]) j = fail[j];
19 if (++j == m) {
20 count++;
21 j = fail[j];
22 }
23 }
24 return count;
25 }
26 }
\end{Code}
\subsection{7.2 AhoCorasick}
\begin{NoteBlock}
Tìm nhiều mẫu ps trong văn bản t.
\end{NoteBlock}
\begin{Code}
1 class AhoCorasick {
2 int m;
3 Node root;
4 AhoCorasick(char[][] ps) {
5 m = ps.length;
6 root = new Node();
7 for (int i = 0; i < m; i++) {
8 Node t = root;
9 for (char c : ps[i]) {
10 if (!t.cs.containsKey(c)) t.cs.put(c, new Node());
11 t = t.cs.get(c);
12 }
13 t.accept.add(i);
14 }
15 Queue<Node> que = new LinkedList<Node>();
16 que.offer(root);
17 while (!que.isEmpty()) {
18 Node t = que.poll();
19 for (Map.Entry<Character, Node> e : t.cs.entrySet()) {
20 char c = e.getKey();
21 Node u = e.getValue();
22 que.offer(u);
23 Node r = t.fail;
24 while (r != null && !r.cs.containsKey(c)) r = r.fail;
25 if (r == null) u.fail = root;
26 else u.fail = r.cs.get(c);
27 u.accept.addAll(u.fail.accept);
28 }
29 }
30 }
31 int[] searchFrom(char[] t) {
32 int n = t.length;
33 int[] count = new int[m];
34 Node u = root;
35 for (int i = 0; i < n; i++) {
36 while (u != null && !u.cs.containsKey(t[i])) u = u.fail;
37 if (u == null) u = root;
38 else u = u.cs.get(t[i]);
39 for (int j : u.accept) count[j]++;
40 }
41 return count;
42 }
43 class Node {
44 Map<Character, Node> cs = new TreeMap<Character, Node>();
45 List<Integer> accept = new ArrayList<Integer>();
46 Node fail;
47 }
48 }
\end{Code}
\subsection{7.3 SuffixArray}
\begin{NoteBlock}
O(n log n)
\end{NoteBlock}
\begin{Code}
1 class SuffixArray {
2 int n;
3 char[] cs;
4 int[] si, is;//ソート後と元のインデックスの対応
5 SuffixArray(char[] t) {
6 n = t.length;
7 cs = new char[n + 1];
8 System.arraycopy(t, 0, cs, 0, n);
9 cs[n] = 0;//番兵 (他の全ての値より小さくすること)
10 is = new int[n + 1];
11 for (int i = 0; i <= n; i++) is[i] = cs[i];
12 si = indexSort(is);
13 int[] a = new int[n + 1], b = new int[n + 1];
14 for (int h = 0; ; ) {
15 for (int i = 0; i < n; i++) {
16 int x = si[i + 1], y = si[i];
17 b[i + 1] = b[i];
18 if (is[x] > is[y] || is[x + h] > is[y + h]) b[i + 1]++;
19 }
20 for (int i = 0; i <= n; i++) is[si[i]] = b[i];
21 if (b[n] == n) break;
22 h = max(1, h << 1);
23 for (int k = h; k >= 0; k -= h) {
24 fill(b, 0);
25 b[0] = k;
26 for (int i = k; i <= n; i++) b[is[i]]++;
27 for (int i = 0; i < n; i++) b[i + 1] += b[i];
28 for (int i = n; i >= 0; i--) {
29 a[--b[si[i] + k > n ? 0 : is[si[i] + k]]] = si[i];
30 }
31 int[] tmp = si; si = a; a = tmp;
32 }
33 }
34 }
35 int[] hs;//hs[i]:=ソート後の suffix の i と i+1 の LCP
36 void buildHs() {
37 hs = new int[n + 1];
38 for (int i = 0, h = 0; i < n; i++) {
39 for (int j = si[is[i] - 1]; cs[i + h] == cs[j + h]; h++);
40 hs[is[i] - 1] = h;
41 if (h > 0) h--;
42 }
43 }
44 RMQ rmq;
45 void buildRMQ() {
46 rmq = new RMQ(hs);//値を返す版
47 }
48 //位置 i,j から始まる部分列の LCP を求める
49 int getLCP(int i, int j) {
50 if (i == j) return n - i;
51 return rmq.query(min(is[i], is[j]), max(is[i], is[j]));
52 }
53 }
\end{Code}
\subsection{7.4 Palindrome}
\begin{NoteBlock}
Tính chất palindrome: nếu [a,b), [a,c), [b,d) đều là palindrome thì [c,d) cũng là palindrome.
Tính độ dài palindrome trong O(n).
len[i]: độ dài palindrome có tâm tại vị trí i/2.
\end{NoteBlock}
\begin{Code}
1 int[] palindrome(char[] cs) {
2 int n = cs.length;
3 int[] len = new int[n * 2];
4 for (int i = 0, j = 0, k; i < n * 2; i += k, j = max(j - k, 0)) {
5 while (i - j >= 0 && i + j + 1 < n * 2
6 && cs[(i - j) / 2] == cs[(i + j + 1) / 2]) j++;
7 len[i] = j;
8 for (k = 1; i - k >= 0 && j - k >= 0 && len[i - k] != j - k; k++) {
9 len[i + k] = min(len[i - k], j - k);
10 }
11 }
12 return len;
13 }
\end{Code}
\section{8 Khác}
\subsection{8.1 Matroid}
\subsubsection{8.1.1 Hệ tiên đề độc lập}
\begin{NoteBlock}
(M1) ∅ ∈ F
(M2) X ⊆ Y ∈ F ⇒ X ∈ F
(M3) X, Y ∈ F, |X| > |Y | ⇒ ∃x ∈ X \ Y. Y ∪ {x} ∈ F
(M3′
Ghi chú/công thức: )   X ⊆ E
Mọi cơ sở của X đều có cùng số phần tử.
\end{NoteBlock}
\subsubsection{8.1.2 Hệ tiên đề hạng}
\begin{NoteBlock}
(R1) r(X) ≤ |X|
(R2) X ⊆ Y ⇒ r(X) ≤ r(Y )
(R3) r(X ∪ Y ) + r(X ∩ Y ) ≤ r(X) + r(Y )
\end{NoteBlock}
\subsubsection{8.1.3 Giao matroid}
\begin{NoteBlock}
Giao của hai matroid có thể không phải matroid, nhưng thuật toán sau tìm được tập độc lập lớn nhất.
Tìm được tập độc lập.
1. Khởi tạo X = empty.
2. Lấy E union {s,t} làm tập đỉnh; dựng cạnh như sau rồi tìm đường ngắn nhất từ s đến t.
(s,y): thêm y vào X vẫn độc lập trong F1.
(y,t): thêm y vào X vẫn độc lập trong F2.
(x,y): thêm y vào X làm mất độc lập trong F1, nhưng
bỏ x thì độc lập trở lại.
(y,x): thêm y vào X làm mất độc lập trong F2, nhưng
bỏ x thì độc lập trở lại.
3. Nếu không tồn tại đường ngắn nhất thì X là tối đại/lớn nhất.
Nếu có, cập nhật X := X union {y0,...,ym} minus {x1,...,xm}, rồi quay lại bước 2.
\end{NoteBlock}
\subsubsection{8.1.4 Phân hoạch matroid}
\begin{NoteBlock}
Ghi chú/công thức: X = X1 ∪ · · · ∪ Xk  Xi ∈ Fi  X   ⇔ X ∈ F
Ghi chú/công thức: E′
Ghi chú/công thức: = E × {1, · · · , k}   2
F1 = {Q ⊆ E′
Ghi chú/công thức: :   i   Qi ∈ Fi}
F2 = {Q ⊆ E′
Ghi chú/công thức: :   i ̸= j   Qi ∩ Qj = ∅}
Ghi chú/công thức: Qi = {e ∈ E′
: (e, i) ∈ Q}
\end{NoteBlock}
\subsection{8.2 Lý thuyết trò chơi tổ hợp}
\subsubsection{8.2.1 Xác định trạng thái thắng}
\begin{NoteBlock}
Thắng khi và chỉ khi tồn tại nước đi đến trạng thái thua.
Thua khi và chỉ khi mọi nước đi đều dẫn đến trạng thái thắng.
\end{NoteBlock}
\subsubsection{8.2.2 Nim}
\begin{NoteBlock}
Trò chơi có nhiều đống; mỗi lượt lấy tùy ý từ một đống, người không lấy được thì thua.
Lấy XOR kích thước các đống; bằng 0 thì thua, khác 0 thì thắng.
Nước tối ưu là chọn đống có bit cao nhất của XOR bật, rồi lấy sao cho XOR trở thành 0.
Ghi chú/công thức: XOR   0   0,1  XOR   0
Ghi chú/công thức: 0,1
\end{NoteBlock}
\subsubsection{8.2.3 Số Grundy}
\begin{NoteBlock}
Trò chơi có số Grundy a tương đương một đống Nim kích thước a.
Số Grundy của trạng thái là mex của các số Grundy của mọi trạng thái đi tới được.
Nếu tách thành nhiều trò chơi, dùng XOR các số Grundy.
Nếu có nhiều trò chơi và mỗi lượt chỉ đi một trò, lấy XOR.
Ghi chú/công thức: 1 - M   2   mod (M + 1)
\end{NoteBlock}
\subsection{8.3 Ngày tháng}
\subsubsection{8.3.1 Đổi sang số ngày}
\begin{NoteBlock}
Đổi sang số ngày, có xét năm nhuận.
\end{NoteBlock}
\begin{Code}
1 int days(int y, int m, int d) {
2 if (m < 3) {
3 y--;
4 m += 12;
5 }
6 return 365 * y + y / 4 - y / 100 + y / 400 + (153 * m + 2) / 5 + d;
7 }
\end{Code}
\subsubsection{8.3.2 java.util.GregorianCalendar}
\begin{NoteBlock}
Thứ trong tuần: [Sunday, Saturday] = [1,7].
Tháng: [January, December] = [0,11].
Sau ngày 4/10/1582 là ngày 15/10/1582; tại đây lịch Julius chuyển sang lịch Gregory.
Ghi chú/công thức: setGregorianChange(new Date(Long.MIN_VALUE))
Ghi chú/công thức: setGregorianChange(new Date(Long.MAX_VALUE))
Ghi chú/công thức: 
Ghi chú/công thức: ( ) 2008   6   31   = 2008   7   1  , 2008   7   31   - 1   = 2008   6   30
Ghi chú/công thức: getActualMaximum(Calendar.DAY_OF_MONTH)
Ghi chú/công thức: isLeapYear(year)
\end{NoteBlock}
\subsection{8.4 Mục nhỏ}
\subsubsection{8.4.1 Scanner}
\begin{NoteBlock}
Chú ý đặc tả của hasNext khác nhau.
\end{NoteBlock}
\begin{Code}
1 class Scanner {
2 BufferedReader br;
3 StringTokenizer st;
4 Scanner(InputStream in) {
5 br = new BufferedReader(new InputStreamReader(in));
6 eat("");
7 }
8 void eat(String s) {
9 st = new StringTokenizer(s);
10 }
11 String nextLine() {
12 try {
13 return br.readLine();
14 } catch (IOException e) {
15 throw new IOError(e);
16 }
17 }
18 boolean hasNext() {
19 while (!st.hasMoreTokens()) {
20 String s = nextLine();
21 if (s == null) return false;
22 eat(s);
23 }
24 return true;
25 }
26 String next() {
27 hasNext();
28 return st.nextToken();
29 }
30 int nextInt() {
31 return Integer.parseInt(next());
32 }
33 long nextLong() {
34 return Long.parseLong(next());
35 }
36 double nextDouble() {
37 return Double.parseDouble(next());
38 }
39 }
\end{Code}
\subsubsection{8.4.2 Liệt kê bit-combination nCk}
\begin{NoteBlock}
Liệt kê theo thứ tự tăng dần.
\end{NoteBlock}
\begin{Code}
1 for (int comb = (1 << k) - 1; comb < 1 << n;) {
2 //...
3 int x = comb & -comb, y = comb + x;
4 comb = ((comb & ~y) / x >>> 1) | y;
5 }
\end{Code}
\subsubsection{8.4.3 Sinh hoán vị}
\begin{NoteBlock}
next permutation của C++.
Sinh hoán vị kế tiếp và trả về liệu nó có tồn tại hay không.
Ngay cả khi có nhiều giá trị bằng nhau, vẫn sinh không trùng.
Muốn sinh toàn bộ hoán vị thì hãy sắp xếp ban đầu.
\end{NoteBlock}
\begin{Code}
1 boolean nextPermutation(int[] is) {
2 int n = is.length;
3 for (int i = n - 1; i > 0; i--) {
4 if (is[i - 1] < is[i]) {
5 int j = n;
6 while (is[i - 1] >= is[--j]);
7 swap(is, i - 1, j);
8 rev(is, i, n);
9 return true;
10 }
11 }
12 rev(is, 0, n);
13 return false;
14 }
\end{Code}
\subsubsection{8.4.4 Xúc xắc}
\begin{NoteBlock}
Sinh toàn bộ trạng thái quay của khối lập phương.
\end{NoteBlock}
\begin{Code}
1 class Dice {
2 //top, front, left, right, back, bottom
3 int[] is = new int[6];
4 Dice(int...is) {
5 this.is = is;
6 }
7 void rollX() {
8 roll(0, 1, 5, 4);
9 }
10 void rollY() {
11 roll(0, 3, 5, 2);
12 }
13 void rollZ() {
14 roll(1, 3, 4, 2);
15 }
16 void roll(int a, int b, int c, int d) {
17 int t = is[d];
18 is[d] = is[c];
19 is[c] = is[b];
20 is[b] = is[a];
21 is[a] = t;
22 }
23 Dice[] rollAll() {
24 Dice[] dices = new Dice[24];
25 for (int i = 0; i < 6; i++) {
26 for (int j = 0; j < 4; j++) {
27 dices[i * 4 + j] = new Dice(is.clone());
28 rollZ();
29 }
30 if ((i & 1) > 0) rollX();
31 else rollY();
32 }
33 return dices;
34 }
35 }
\end{Code}
\subsubsection{8.4.5 Số Josephus}
\begin{NoteBlock}
Ghi chú/công thức: n   m   k   [0, n - 1]
J(n, m, 0) = -1
J(n, m, k) = (J(n - 1, m, k - 1) + k)%n
O(k)
\end{NoteBlock}
\begin{Code}
1 int josephus(int n, int m, int k) {
2 int x = -1;
3 for (int i = n - k + 1; i <= n; i++) {
4 x = (x + m) % i;
5 }
6 return x;
7 }
\end{Code}
\begin{NoteBlock}
Số Josephus ngược.
O(n)
Ghi chú/công thức: n   m   x   [1, n]
\end{NoteBlock}
\begin{Code}
1 int invJosephus(int n, int m, int x) {
2 for (int i = n; ; i--) {
3 if (x == i) return n - i;
4 x = (x - m % i + i) % i;
5 }
6 }
\end{Code}
\subsubsection{8.4.6 Hình chữ nhật lớn nhất}
\begin{NoteBlock}
Cho mảng chiều cao; tìm diện tích hình chữ nhật lớn nhất chứa trong đó.
O(n)
\end{NoteBlock}
\begin{Code}
1 int maxRect(int[] _hs) {
2 int n = _hs.length;
3 int res = 0;
4 int[] hs = new int[n + 1];//n 番に番兵 0 を置く
5 System.arraycopy(_hs, 0, hs, 0, n);
6 Stack<Integer> sx = new Stack<Integer>();
7 Stack<Integer> sh = new Stack<Integer>();
8 sh.push(0);
9 for (int i = 0; i <= n; i++) {
10 int x = i;
11 while (sh.peek() > hs[i]) {
12 res = max(res, sh.pop() * (i - (x = sx.pop())));
13 }
14 sx.push(x);
15 sh.push(hs[i]);
16 }
17 return res;
18 }
\end{Code}
\section{9 Tìm kiếm}
\subsection{9.1 Tìm kiếm alpha-beta}
\begin{Code}
1 int alphaBeta(State s, int alpha, int beta) {
2 if (s.finished()) return s.calcScore();
3 for (State t : s.next()) {
4 alpha = max(alpha, -alphaBeta(t, -beta, -alpha));
5 if (alpha >= beta) break;
6 }
7 return alpha;
8 }
\end{Code}
\subsection{9.2 Bài toán tô màu đỉnh tối thiểu}
\begin{NoteBlock}
Tô màu bằng số màu nhỏ nhất sao cho hai đỉnh kề nhau không cùng màu.
Ý tưởng tìm kiếm.
1. Tách thành các thành phần liên thông rồi giải.
2. Sắp theo thứ tự MA rồi bắt đầu tìm kiếm.
3. DFS có cắt tỉa theo số màu.
\end{NoteBlock}
\begin{Code}
1 int[] coloring(boolean[][] g) {
2 this.g = g;
3 int n = g.length;
4 res = new int[n];
5 id = new int[n + 1];//元のインデックス (n 番は番兵)
6 int[] b = new int[n + 1];//連結度
7 for (int i = 0; i <= n; i++) id[i] = i;
8 for (s = 0, t = 1; t <= n; t++) {//MA 順序でソート
9 int p = t;//最大連結度の点
10 for (int i = t; i < n; i++) {
11 if (g[id[t - 1]][id[i]]) b[id[i]]++;
12 if (b[id[p]] < b[id[i]]) p = i;
13 }
14 swap(id, t, p);
15 if (b[id[t]] == 0) {//連結度 0 なので [s,t) が連結成分
16 min = n + 1;
17 dfs(new int[n], s, 0);
18 s = t;
19 }
20 }
21 return res;
22 }
23 void dfs(int[] is, int p, int k) {//is:現在の塗り方,p:次塗る点,k:使った色の数
24 if (k >= min) return;
25 if (p == t) {
26 for (int i = s; i < t; i++) res[id[i]] = is[i];
27 min = k;
28 } else {
29 boolean[] used = new boolean[k + 1];
30 for (int i = 0; i < p; i++) if (g[id[p]][id[i]]) used[is[i]] = true;
31 for (int i = 0; i <= k; i++) if (!used[i]) {//今までに使った色 +1 まで調べる
32 int[] js = is.clone();
33 js[p] = i;
34 dfs(js, p + 1, max(k, i + 1));
35 }
36 }
37 }
\end{Code}
\subsection{9.3 Số matching hoàn hảo}
\begin{NoteBlock}
Giải tổng quát các bài DP bit như đếm matching hoàn hảo hoặc đếm tập độc lập.
\end{NoteBlock}
\begin{Code}
1 long match(V[] vs) {
2 int n = vs.length;
3 for (int i = 0; i < n; i++) vs[i].id = i;
4 for (V v : vs) for (V u : v) v.max = max(v.max, u.id);
5 long[] crt = {1};
6 int[] cvi = new int[n], civ = new int[n];
7 int cn = 0;
8 fill(cvi, -1);
9 for (int k = 0; k < n; k++) {
10 int nn = 0;
11 int[] nvi = new int[n], niv = new int[n];
12 fill(nvi, -1);
13 for (int i = 0; i <= k; i++) if (vs[i].max > k) {
14 niv[nn] = i;
15 nvi[i] = nn++;
16 }
17 long[] next = new long[1 << nn];
18 loop : for (int i = 0; i < 1 << cn; i++) if (crt[i] > 0) {
19 int i2 = 0;
20 boolean need = false;
21 for (int j = 0; j < cn; j++) if ((i >> j & 1) != 0) {
22 if (nvi[civ[j]] < 0) {
23 if (need) continue loop;
24 need = true;
25 } else {
26 i2 |= 1 << nvi[civ[j]];
27 }
28 }
29 if (need) {
30 next[i2] += crt[i];
31 } else {
32 if (nvi[k] >= 0) {
33 next[i2 | 1 << nvi[k]] += crt[i];
34 }
35 for (V u : vs[k]) if (u.id < k && (i >> cvi[u.id] & 1) != 0) {
36 next[i2 ^ 1 << nvi[u.id]] += crt[i];
37 }
38 }
39 }
40 crt = next; cvi = nvi; civ = niv; cn = nn;
41 }
42 return crt[0];
43 }
44 class V extends ArrayList<V> {
45 int id, max = -1;
46 }
\end{Code}
\section{10 Trình trực quan hóa}
\begin{Code}
1 class Vis extends JFrame {
2 BufferedImage img;
3 Graphics2D g;
4 boolean stop;
5 Vis() {
6 img = new BufferedImage(500, 500, BufferedImage.TYPE_INT_RGB);
7 g = (Graphics2D)img.getGraphics();
8 g.setRenderingHint(RenderingHints.KEY_ANTIALIASING,
9 RenderingHints.VALUE_ANTIALIAS_ON);
10 clear();
11 JComponent c = new JComponent() {
12 protected void paintComponent(Graphics g) {
13 super.paintComponent(g);
14 g.drawImage(img, 0, 0, null);
15 }
16 };
17 c.setPreferredSize(new Dimension(500, 500));
18 add(c);
19 addMouseListener(new MouseAdapter() {
20 public void mouseClicked(MouseEvent e) {
21 stop = false;
22 }
23 });
24 setDefaultCloseOperation(EXIT_ON_CLOSE);
25 pack();
26 setVisible(true);
27 }
28 void vis() {
29 repaint();
30 stop = true;
31 try {
32 while (stop) Thread.sleep(10);
33 } catch (Exception e) {
34 }
35 }
36 void clear() {
37 g.setColor(Color.WHITE);
38 g.fillRect(0, 0, 500, 500);
39 g.setColor(Color.BLACK);
40 }
41 }
\end{Code}
\section{11 Bảng tiện dụng}
\subsection{11.1 Định dạng}
\begin{NoteBlock}
'-': căn trái.
'+': luôn in dấu.
' ': thêm khoảng trắng trước số không âm.
'0': điền 0 ở bên trái.
',': phân tách mỗi 1000.
\end{NoteBlock}
\subsection{11.2 Biểu thức chính quy}
\subsubsection{11.2.1 Lớp ký tự}
\begin{NoteBlock}
[abc]: a, b hoặc c (lớp đơn giản).
[^abc]: ký tự không phải a, b, c (phủ định).
Ghi chú/công thức: [a-zA-Z] a - z   A - Z ( )
Ghi chú/công thức: [a-d[m-p]] a - d  m - p:[a-dm-p] ( )
[a-z&&[def]]: d, e, f (giao).
Ghi chú/công thức: [a-z&&[^bc]] b   c   a - z:[ad-z] ( )
Ghi chú/công thức: [a-z&&[^m-p]] m - p   a - z:[a-lq-z] ( )
.: ký tự bất kỳ.
\d: chữ số [0-9].
\D: không phải chữ số [^0-9].
\s: ký tự trắng [ \t\n\x0B\f\r].
\S: không phải ký tự trắng [^\s].
\w: ký tự từ [a-zA-Z_0-9].
\W: không phải ký tự từ [^\w].
\end{NoteBlock}
\subsubsection{11.2.2 Biên}
\begin{NoteBlock}
^: đầu dòng.
$: cuối dòng.
\b: biên từ.
\B: không phải biên từ.
\A: đầu vào.
\G: cuối lần match trước.
\Z: cuối đầu vào, trừ ký hiệu xuống dòng cuối nếu có.
\z: cuối đầu vào.
\end{NoteBlock}
\subsubsection{11.2.3 Lượng từ}
\begin{NoteBlock}
Không chỉ định gì thì match dài nhất.
Thêm ? phía sau thì match ngắn nhất.
Thêm + phía sau thì trở thành possessive/greedy cứng, không quay lui khi match thất bại.
X?: X xuất hiện 1 hoặc 0 lần.
X*: X xuất hiện 0 lần trở lên.
X+: X xuất hiện 1 lần trở lên.
X{n}: X xuất hiện n lần.
X{n,}: X xuất hiện ít nhất n lần.
X{n,m}: X xuất hiện từ n đến m lần.
\end{NoteBlock}
\subsubsection{11.2.4 Khác}
\begin{NoteBlock}
Nhóm capture số 0 là toàn bộ match; các nhóm được đánh số từ 1 theo thứ tự mở ngoặc.
XY: Y ngay sau X.
X|Y: X hoặc Y.
(X): X, nhóm capture.
(?:X): X, nhóm không capture.
\n: nhóm capture thứ n đã match.
(?i): bật CASE_INSENSITIVE.
(?-i): tắt CASE_INSENSITIVE.
\end{NoteBlock}
\subsection{11.3 Bảng số}
\begin{NoteBlock}
n log10 n n! nC(n/2) LCM(1 . . . n) Pn Bn
\end{NoteBlock}
\begin{Code}
2 0.30 2 2 2 2 2
3 0.48 6 3 6 3 5
4 0.60 24 6 12 5 15
5 0.70 120 10 60 7 52
6 0.78 720 20 60 11 203
7 0.85 5040 35 420 15 877
8 0.90 40320 70 840 22 4140
9 0.95 362880 126 2520 30 21147
10 1.00 3628800 252 2520 42 115975
11 39916800 462 27720 56 678570
12 479001600 924 27720 77 4213597
15 6435 360360 176 1382958545
20 184756 232792560 627
25 5200300 1958
30 155117520 5604
40 37338
50 204226
70 4087968
100 190569292
\end{Code}
\end{document}
